
\Data\ID{1003021103}
\Updated{2020-07-15 03:07:12}
\Level{A}
\SubArea{Limits and continuity}
\LANGen{
\Question{
Evaluate the following limit.
\[ \lim\limits_{x\rightarrow-1}\frac{x^2+x-2}{x^3-1} \]}
{\( 1 \)}{\( 0 \)}{\( 2 \)}{\( -2 \)}{}{}}
\LANGpl{
\Question{
Oblicz granicę funkcji.
\[ \lim\limits_{x\rightarrow-1}\frac{x^2+x-2}{x^3-1} \]}
{\( 1 \)}{\( 0 \)}{\( 2 \)}{\( -2 \)}{}{}}
\LANGcs{
\Question{
Vypočítejte následující limitu.
\[ \lim\limits_{x\rightarrow-1}\frac{x^2+x-2}{x^3-1} \]}
{\( 1 \)}{\( 0 \)}{\( 2 \)}{\( -2 \)}{}{}}
\LANGsk{
\Question{
Vypočítajte následujúcu limitu
\[ \lim\limits_{x\rightarrow-1}\frac{x^2+x-2}{x^3-1} \]}
{\( 1 \)}{\( 0 \)}{\( 2 \)}{\( -2 \)}{}{}}
\LANGes{
\Question{
Resuelve el siguiente límite:
\[ \lim\limits_{x\rightarrow-1}\frac{x^2+x-2}{x^3-1} \]}
{\( 1 \)}{\( 0 \)}{\( 2 \)}{\( -2 \)}{}{}}
\End

\Data\ID{1003021104}
\Updated{2020-07-15 03:07:12}
\Level{A}
\SubArea{Limits and continuity}
\LANGen{
\Question{
Evaluate the following limit.
\[ \lim\limits_{x\rightarrow-2}\frac{3x^2+3x-6}{2x^2-2x-12} \]}
{\( \frac9{10} \)}{\( -\frac32 \)}{\( \frac32 \)}{\( 0 \)}{}{}}
\LANGpl{
\Question{
Oblicz granicę funkcji.
\[ \lim\limits_{x\rightarrow-2}\frac{3x^2+3x-6}{2x^2-2x-12} \]}
{\( \frac9{10} \)}{\( -\frac32 \)}{\( \frac32 \)}{\( 0 \)}{}{}}
\LANGcs{
\Question{
Vypočítejte následující limitu.
\[ \lim\limits_{x\rightarrow-2}\frac{3x^2+3x-6}{2x^2-2x-12} \]}
{\( \frac9{10} \)}{\( -\frac32 \)}{\( \frac32 \)}{\( 0 \)}{}{}}
\LANGsk{
\Question{
Vypočítajte následujúcu limitu
\[ \lim\limits_{x\rightarrow-2}\frac{3x^2+3x-6}{2x^2-2x-12} \]}
{\( \frac9{10} \)}{\( -\frac32 \)}{\( \frac32 \)}{\( 0 \)}{}{}}
\LANGes{
\Question{
Resuelve el siguiente límite:
\[ \lim\limits_{x\rightarrow-2}\frac{3x^2+3x-6}{2x^2-2x-12} \]}
{\( \frac9{10} \)}{\( -\frac32 \)}{\( \frac32 \)}{\( 0 \)}{}{}}
\End

\Data\ID{1003021105}
\Updated{2020-07-15 03:07:12}
\Level{A}
\SubArea{Limits and continuity}
\LANGen{
\Question{
Evaluate the following limit.
\[ \lim\limits_{x\rightarrow1}\frac{x^2+x-2}{1-x^3} \]}
{\( -1 \)}{\( 1 \)}{\( 0 \)}{\( 3 \)}{}{}}
\LANGpl{
\Question{
Oblicz granicę funkcji.
\[ \lim\limits_{x\rightarrow1}\frac{x^2+x-2}{1-x^3} \]}
{\( -1 \)}{\( 1 \)}{\( 0 \)}{\( 3 \)}{}{}}
\LANGcs{
\Question{
Vypočítejte následující limitu.
\[ \lim\limits_{x\rightarrow1}\frac{x^2+x-2}{1-x^3} \]}
{\( -1 \)}{\( 1 \)}{\( 0 \)}{\( 3 \)}{}{}}
\LANGsk{
\Question{
Vypočítajte následujúcu limitu.
\[ \lim\limits_{x\rightarrow1}\frac{x^2+x-2}{1-x^3} \]}
{\( -1 \)}{\( 1 \)}{\( 0 \)}{\( 3 \)}{}{}}
\LANGes{
\Question{
Resuelve el siguiente límite:
\[ \lim\limits_{x\rightarrow1}\frac{x^2+x-2}{1-x^3} \]}
{\( -1 \)}{\( 1 \)}{\( 0 \)}{\( 3 \)}{}{}}
\End

\Data\ID{1003021106}
\Updated{2020-07-15 03:07:12}
\Level{A}
\SubArea{Limits and continuity}
\LANGen{
\Question{
Evaluate the following limit.
\[ \lim\limits_{x\rightarrow-3}\frac{x^2+5x+6}{x^3+27} \]}
{\( -\frac1{27} \)}{\( -\frac5{27} \)}{\( -\frac19 \)}{\( \frac1{27} \)}{}{}}
\LANGpl{
\Question{
Oblicz granicę funkcji.
\[ \lim\limits_{x\rightarrow-3}\frac{x^2+5x+6}{x^3+27} \]}
{\( -\frac1{27} \)}{\( -\frac5{27} \)}{\( -\frac19 \)}{\( \frac1{27} \)}{}{}}
\LANGcs{
\Question{
Vypočítejte následující limitu.
\[ \lim\limits_{x\rightarrow-3}\frac{x^2+5x+6}{x^3+27} \]}
{\( -\frac1{27} \)}{\( -\frac5{27} \)}{\( -\frac19 \)}{\( \frac1{27} \)}{}{}}
\LANGsk{
\Question{
Vypočítajte následujúcu limitu.
\[ \lim\limits_{x\rightarrow-3}\frac{x^2+5x+6}{x^3+27} \]}
{\( -\frac1{27} \)}{\( -\frac5{27} \)}{\( -\frac19 \)}{\( \frac1{27} \)}{}{}}
\LANGes{
\Question{
Resuelve el siguiente límite:
\[ \lim\limits_{x\rightarrow-3}\frac{x^2+5x+6}{x^3+27} \]}
{\( -\frac1{27} \)}{\( -\frac5{27} \)}{\( -\frac19 \)}{\( \frac1{27} \)}{}{}}
\End

\Data\ID{1003093101}
\Updated{2021-05-17 08:05:53}
\Level{A}
\SubArea{Limits and continuity}
\LANGen{
\Question{
Which of the statements A, B, C, D, E given bellow are correct?
\[ \begin{aligned}
\text{A: }&  \lim\limits_{x\to1^+}\frac{2x}{x-1}=\infty  \\
\text{B: }&  \lim\limits_{x\to-2^+}\frac{x+1}{x+2}=-\infty \\
\text{C: }&  \lim\limits_{x\to3^-}\frac{x+1}{3-x}=-\infty \\
\text{D: }&  \lim\limits_{x\to-1^+}\frac{x}{x^2-1}=\infty \\
\text{E: }&  \lim\limits_{x\to2^-}\frac{x+3}{4-x^2} =\infty 
\end{aligned} \]
The only correct statements are:}
{A, B, D, E}{A, B, D}{A, D, E}{A, B, C, E}{A, B, C}{B, D, E}}
\LANGpl{
\Question{
Które z podanych wyrażeń  A, B, C, D, E są poprawne?
\[ \begin{aligned}
\text{A: }&  \lim\limits_{x\to1^+}\frac{2x}{x-1}=\infty  \\
\text{B: }&  \lim\limits_{x\to-2^+}\frac{x+1}{x+2}=-\infty \\
\text{C: }&  \lim\limits_{x\to3^-}\frac{x+1}{3-x}=-\infty \\
\text{D: }&  \lim\limits_{x\to-1^+}\frac{x}{x^2-1}=\infty \\
\text{E: }&  \lim\limits_{x\to2^-}\frac{x+3}{4-x^2} =\infty 
\end{aligned} \]
Poprawnymi wyrażeniami są:}
{A, B, D, E}{A, B, D}{A, D, E}{A, B, C, E}{A, B, C}{B, D, E}}
\LANGcs{
\Question{
Která z níže uvedených tvrzení A, B, C, D, E jsou pravdivá?
\[ \begin{aligned}
\text{A: }&  \lim\limits_{x\to1^+}\frac{2x}{x-1}=\infty  \\
\text{B: }&  \lim\limits_{x\to-2^+}\frac{x+1}{x+2}=-\infty \\
\text{C: }&  \lim\limits_{x\to3^-}\frac{x+1}{3-x}=-\infty \\
\text{D: }&  \lim\limits_{x\to-1^+}\frac{x}{x^2-1}=\infty \\
\text{E: }&  \lim\limits_{x\to2^-}\frac{x+3}{4-x^2} =\infty 
\end{aligned} \]
Jedinými pravdivými tvrzeními jsou:}
{A, B, D, E}{A, B, D}{A, D, E}{A, B, C, E}{A, B, C}{B, D, E}}
\LANGsk{
\Question{
Ktoré z tvrdení A, B, C, D, E uvedených nižšie sú správne?
\[ \begin{aligned}
\text{A: }&  \lim\limits_{x\to1^+}\frac{2x}{x-1}=\infty  \\
\text{B: }&  \lim\limits_{x\to-2^+}\frac{x+1}{x+2}=-\infty \\
\text{C: }&  \lim\limits_{x\to3^-}\frac{x+1}{3-x}=-\infty \\
\text{D: }&  \lim\limits_{x\to-1^+}\frac{x}{x^2-1}=\infty \\
\text{E: }&  \lim\limits_{x\to2^-}\frac{x+3}{4-x^2} =\infty 
\end{aligned} \]
Jedinými správnymi tvrdeniami sú:}
{A, B, D, E}{A, B, D}{A, D, E}{A, B, C, E}{A, B, C}{B, D, E}}
\LANGes{
\Question{
¿Cuáles de las siguientes igualdades A, B, C, D, E son correctas?
\[ \begin{aligned}
\text{A: }&  \lim\limits_{x\to1^+}\frac{2x}{x-1}=\infty  \\
\text{B: }&  \lim\limits_{x\to-2^+}\frac{x+1}{x+2}=-\infty \\
\text{C: }&  \lim\limits_{x\to3^-}\frac{x+1}{3-x}=-\infty \\
\text{D: }&  \lim\limits_{x\to-1^+}\frac{x}{x^2-1}=\infty \\
\text{E: }&  \lim\limits_{x\to2^-}\frac{x+3}{4-x^2} =\infty 
\end{aligned} \]}
{A, B, D, E}{A, B, D}{A, D, E}{A, B, C, E}{A, B, C}{B, D, E}}
\End

\Data\ID{1003093103}
\Updated{2020-07-15 03:07:21}
\Level{A}
\SubArea{Limits and continuity}
\LANGen{
\Question{
Evaluate the following limit.
\[ \lim\limits_{x\to\infty}\left(-2x^3+5x^2-7\right) \]}
{\( -\infty \)}{\( \infty \)}{\( -7 \)}{\( 0 \)}{}{}}
\LANGpl{
\Question{
Oblicz granicę funkcji.
\[ \lim\limits_{x\to\infty}\left(-2x^3+5x^2-7\right) \]}
{\( -\infty \)}{\( \infty \)}{\( -7 \)}{\( 0 \)}{}{}}
\LANGcs{
\Question{
Vypočítejte následující limitu.
\[ \lim\limits_{x\to\infty}\left(-2x^3+5x^2-7\right) \]}
{\( -\infty \)}{\( \infty \)}{\( -7 \)}{\( 0 \)}{}{}}
\LANGsk{
\Question{
Vypočítajte následujúcu limitu.
\[ \lim\limits_{x\to\infty}\left(-2x^3+5x^2-7\right) \]}
{\( -\infty \)}{\( \infty \)}{\( -7 \)}{\( 0 \)}{}{}}
\LANGes{
\Question{
Resuelve el siguiente límite:
\[ \lim\limits_{x\to\infty}\left(-2x^3+5x^2-7\right) \]}
{\( -\infty \)}{\( \infty \)}{\( -7 \)}{\( 0 \)}{}{}}
\End

\Data\ID{1003093104}
\Updated{2020-07-15 03:07:21}
\Level{A}
\SubArea{Limits and continuity}
\LANGen{
\Question{
Evaluate the following limit.
\[ \lim\limits_{x\to-\infty}\left(x^4-3x^2-x+5\right) \]}
{\( \infty \)}{\( -\infty \)}{\( 5 \)}{\( 0 \)}{}{}}
\LANGpl{
\Question{
Oblicz granicę funkcji.
\[ \lim\limits_{x\to-\infty}\left(x^4-3x^2-x+5\right) \]}
{\( \infty \)}{\( -\infty \)}{\( 5 \)}{\( 0 \)}{}{}}
\LANGcs{
\Question{
Vypočítejte následující limitu.
\[ \lim\limits_{x\to-\infty}\left(x^4-3x^2-x+5\right) \]}
{\( \infty \)}{\( -\infty \)}{\( 5 \)}{\( 0 \)}{}{}}
\LANGsk{
\Question{
Vypočítajte následujúcu limitu.
\[ \lim\limits_{x\to-\infty}\left(x^4-3x^2-x+5\right) \]}
{\( \infty \)}{\( -\infty \)}{\( 5 \)}{\( 0 \)}{}{}}
\LANGes{
\Question{
Resuelve el siguiente límite:
\[ \lim\limits_{x\to-\infty}\left(x^4-3x^2-x+5\right) \]}
{\( \infty \)}{\( -\infty \)}{\( 5 \)}{\( 0 \)}{}{}}
\End

\Data\ID{1003093105}
\Updated{2020-07-15 03:07:21}
\Level{A}
\SubArea{Limits and continuity}
\LANGen{
\Question{
Evaluate the following limit.
\[ \lim\limits_{x\to\infty}\frac{3x^2-1}{2-x} \]}
{\( -\infty \)}{\( \infty \)}{\( 3 \)}{\( 0 \)}{}{}}
\LANGpl{
\Question{
Oblicz granicę funkcji.
\[ \lim\limits_{x\to\infty}\frac{3x^2-1}{2-x} \]}
{\( -\infty \)}{\( \infty \)}{\( 3 \)}{\( 0 \)}{}{}}
\LANGcs{
\Question{
Vypočítejte následující limitu.
\[ \lim\limits_{x\to\infty}\frac{3x^2-1}{2-x} \]}
{\( -\infty \)}{\( \infty \)}{\( 3 \)}{\( 0 \)}{}{}}
\LANGsk{
\Question{
Vypočítajte následujúcu limitu.
\[ \lim\limits_{x\to\infty}\frac{3x^2-1}{2-x} \]}
{\( -\infty \)}{\( \infty \)}{\( 3 \)}{\( 0 \)}{}{}}
\LANGes{
\Question{
Resuelve el siguiente límite:
\[ \lim\limits_{x\to\infty}\frac{3x^2-1}{2-x} \]}
{\( -\infty \)}{\( \infty \)}{\( 3 \)}{\( 0 \)}{}{}}
\End

\Data\ID{1003093106}
\Updated{2020-07-15 03:07:21}
\Level{A}
\SubArea{Limits and continuity}
\LANGen{
\Question{
Evaluate the following limit.
\[ \lim\limits_{x\to-\infty}\frac{3x^2+x-5}{x^2-2x+1} \]}
{\( 3 \)}{\( \infty \)}{\( -\infty \)}{\( 0 \)}{}{}}
\LANGpl{
\Question{
Oblicz granicę funkcji.
\[ \lim\limits_{x\to-\infty}\frac{3x^2+x-5}{x^2-2x+1} \]}
{\( 3 \)}{\( \infty \)}{\( -\infty \)}{\( 0 \)}{}{}}
\LANGcs{
\Question{
Vypočítejte následující limitu.
\[ \lim\limits_{x\to-\infty}\frac{3x^2+x-5}{x^2-2x+1} \]}
{\( 3 \)}{\( \infty \)}{\( -\infty \)}{\( 0 \)}{}{}}
\LANGsk{
\Question{
Vypočítajte následujúcu limitu.
\[ \lim\limits_{x\to-\infty}\frac{3x^2+x-5}{x^2-2x+1} \]}
{\( 3 \)}{\( \infty \)}{\( -\infty \)}{\( 0 \)}{}{}}
\LANGes{
\Question{
Resuelve el siguiente límite:
\[ \lim\limits_{x\to-\infty}\frac{3x^2+x-5}{x^2-2x+1} \]}
{\( 3 \)}{\( \infty \)}{\( -\infty \)}{\( 0 \)}{}{}}
\End

\Data\ID{1003093107}
\Updated{2020-07-15 03:07:21}
\Level{A}
\SubArea{Limits and continuity}
\LANGen{
\Question{
Evaluate the following limit.
\[ \lim\limits_{x\to-\infty}\frac{x^2-2x+5}{2x^3-x^2+4} \]}
{\( 0 \)}{\( \infty \)}{\( -\infty \)}{\( \frac12 \)}{}{}}
\LANGpl{
\Question{
Oblicz granicę funkcji.
\[ \lim\limits_{x\to-\infty}\frac{x^2-2x+5}{2x^3-x^2+4} \]}
{\( 0 \)}{\( \infty \)}{\( -\infty \)}{\( \frac12 \)}{}{}}
\LANGcs{
\Question{
Vypočítejte následující limitu.
\[ \lim\limits_{x\to-\infty}\frac{x^2-2x+5}{2x^3-x^2+4} \]}
{\( 0 \)}{\( \infty \)}{\( -\infty \)}{\( \frac12 \)}{}{}}
\LANGsk{
\Question{
Vypočítajte následujúcu limitu.
\[ \lim\limits_{x\to-\infty}\frac{x^2-2x+5}{2x^3-x^2+4} \]}
{\( 0 \)}{\( \infty \)}{\( -\infty \)}{\( \frac12 \)}{}{}}
\LANGes{
\Question{
Resuelve el siguiente límite:
\[ \lim\limits_{x\to-\infty}\frac{x^2-2x+5}{2x^3-x^2+4} \]}
{\( 0 \)}{\( \infty \)}{\( -\infty \)}{\( \frac12 \)}{}{}}
\End

\Data\ID{1103024501}
\Updated{2020-07-15 03:07:12}
\Level{A}
\SubArea{Limits and continuity}
\IMG{1103024501.pdf}
\LANGen{
\Question{
Given the graph of the function \( f \), find \( \lim\limits_{x\rightarrow0}f(x) \).
\[f(x)=x^2\cdot\cos\!\left(\frac1x\right)+a,\ a\in\mathbb{R}\]}
{\( a \)}{does not exist}{\( 0 \)}{\( a^2 \)}{}{}}
\LANGpl{
\Question{
Dany jest wykres funkcji \( f \). Wyznacz \( \lim\limits_{x\rightarrow0}f(x) \).
\[f(x)=x^2\cdot\cos\!\left(\frac1x\right)+a,\ a\in\mathbb{R}\]}
{\( a \)}{nie istnieje}{\( 0 \)}{\( a^2 \)}{}{}}
\LANGcs{
\Question{
Je dán graf funkce \( f \). Určete  \( \lim\limits_{x\rightarrow0}f(x) \).
\[f(x)=x^2\cdot\cos\!\left(\frac1x\right)+a,\ a\in\mathbb{R}\]}
{\( a \)}{neexistuje}{\( 0 \)}{\( a^2 \)}{}{}}
\LANGsk{
\Question{
Je daný graf funkcie \( f \). Určte  \( \lim\limits_{x\rightarrow0}f(x) \).
\[f(x)=x^2\cdot\cos\!\left(\frac1x\right)+a,\ a\in\mathbb{R}\]}
{\( a \)}{neexistuje}{\( 0 \)}{\( a^2 \)}{}{}}
\LANGes{
\Question{
Dado el gráfico de la función \( f \), halla \( \lim\limits_{x\rightarrow0}f(x) \).
\[f(x)=x^2\cdot\cos\!\left(\frac1x\right)+a,\ a\in\mathbb{R}\]}
{\( a \)}{no existe.}{\( 0 \)}{\( a^2 \)}{}{}}
\End

\Data\ID{1103024502}
\Updated{2020-07-15 03:07:12}
\Level{A}
\SubArea{Limits and continuity}
\IMG{1103024502.pdf}
\LANGen{
\Question{
Given the graph of the function \( f \), find \( \lim\limits_{x\rightarrow0}f(x) \).
\[f(x)=\sin\!\left(\frac1x\right)\]}
{does not exist}{\( a \)}{\( -a \)}{\( 0 \)}{}{}}
\LANGpl{
\Question{
Dany jest wykres funkcji \( f \). Wyznacz \( \lim\limits_{x\rightarrow0}f(x) \).
\[f(x)=\sin\!\left(\frac1x\right)\]}
{nie istnieje}{\( a \)}{\( -a \)}{\( 0 \)}{}{}}
\LANGcs{
\Question{
Je dán graf funkce \( f \). Určete  \( \lim\limits_{x\rightarrow0}f(x) \).
\[f(x)=\sin\!\left(\frac1x\right)\]}
{neexistuje}{\( a \)}{\( -a \)}{\( 0 \)}{}{}}
\LANGsk{
\Question{
Je daný graf funkcie \( f \). Určte  \( \lim\limits_{x\rightarrow0}f(x) \).
\[f(x)=\sin\!\left(\frac1x\right)\]}
{neexistuje}{\( a \)}{\( -a \)}{\( 0 \)}{}{}}
\LANGes{
\Question{
Dado el gráfico de la función \( f \), halla \( \lim\limits_{x\rightarrow0}f(x) \).
\[f(x)=\sin\!\left(\frac1x\right)\]}
{no existe}{\( a \)}{\( -a \)}{\( 0 \)}{}{}}
\End

\Data\ID{1103024503}
\Updated{2020-07-15 03:07:12}
\Level{A}
\SubArea{Limits and continuity}
\IMG{1103024503.pdf}
\LANGen{
\Question{
Given the graph of the function \( f \), find \( \lim\limits_{x\rightarrow1}f(x) \).}
{does not exist}{\( -1 \)}{\( 0 \)}{\( 2 \)}{}{}}
\LANGpl{
\Question{
Dany jest wykres funkcji \( f \). Wyznacz \( \lim\limits_{x\rightarrow1}f(x) \).}
{nie istnieje}{\( -1 \)}{\( 0 \)}{\( 2 \)}{}{}}
\LANGcs{
\Question{
Je dán graf funkce \( f \). Určete \( \lim\limits_{x\rightarrow1}f(x) \).}
{neexistuje}{\( -1 \)}{\( 0 \)}{\( 2 \)}{}{}}
\LANGsk{
\Question{
Je daný graf funkcie \( f \). Určte \( \lim\limits_{x\rightarrow1}f(x) \).}
{neexistuje}{\( -1 \)}{\( 0 \)}{\( 2 \)}{}{}}
\LANGes{
\Question{
Dado el gráfico de la función \( f \), halla \( \lim\limits_{x\rightarrow1}f(x) \).}
{no existe}{\( -1 \)}{\( 0 \)}{\( 2 \)}{}{}}
\End

\Data\ID{1103024504}
\Updated{2020-07-15 03:07:12}
\Level{A}
\SubArea{Limits and continuity}
\IMG{1103024504.pdf}
\LANGen{
\Question{
Given the graph of the function \( f \), find \( \lim\limits_{x\rightarrow1^-}f(x) \).}
{\( -1 \)}{\( 0 \)}{\( 2 \)}{does not exist}{}{}}
\LANGpl{
\Question{
Dany jest wykres funkcji \( f \). Wyznacz \( \lim\limits_{x\rightarrow1^-}f(x) \).}
{\( -1 \)}{\( 0 \)}{\( 2 \)}{nie istnieje}{}{}}
\LANGcs{
\Question{
Je dán graf funkce \( f \). Určete \( \lim\limits_{x\rightarrow1^-}f(x) \).}
{\( -1 \)}{\( 0 \)}{\( 2 \)}{neexistuje}{}{}}
\LANGsk{
\Question{
Je daný graf funkcie \( f \). Určte \( \lim\limits_{x\rightarrow1^-}f(x) \).}
{\( -1 \)}{\( 0 \)}{\( 2 \)}{neexistuje}{}{}}
\LANGes{
\Question{
Dado el gráfico de la función \( f \), halla \( \lim\limits_{x\rightarrow1^-}f(x) \).}
{\( -1 \)}{\( 0 \)}{\( 2 \)}{no existe}{}{}}
\End

\Data\ID{1103024505}
\Updated{2020-07-15 03:07:12}
\Level{A}
\SubArea{Limits and continuity}
\IMG{1103024505.pdf}
\LANGen{
\Question{
Given the graph of the function \( f \), find \( \lim\limits_{x\rightarrow1^+}f(x) \).}
{\( 2 \)}{\( 0 \)}{\( -1 \)}{does not exist}{}{}}
\LANGpl{
\Question{
Dany jest wykres funkcji \( f \). Wyznacz \( \lim\limits_{x\rightarrow1^+}f(x) \).}
{\( 2 \)}{\( 0 \)}{\( -1 \)}{nie istnieje}{}{}}
\LANGcs{
\Question{
Je dán graf funkce \( f \). Určete \( \lim\limits_{x\rightarrow1^+}f(x) \).}
{\( 2 \)}{\( 0 \)}{\( -1 \)}{neexistuje}{}{}}
\LANGsk{
\Question{
Je daný graf funkcie \( f \). Určte \( \lim\limits_{x\rightarrow1^+}f(x) \).}
{\( 2 \)}{\( 0 \)}{\( -1 \)}{neexistuje}{}{}}
\LANGes{
\Question{
Dado el gráfico de la función \( f \), halla \( \lim\limits_{x\rightarrow1^+}f(x) \).}
{\( 2 \)}{\( 0 \)}{\( -1 \)}{no existe}{}{}}
\End

\Data\ID{1103024506}
\Updated{2020-07-15 03:07:12}
\Level{A}
\SubArea{Limits and continuity}
\IMG{1103024506.pdf}
\LANGen{
\Question{
Given the graph of the function \( f \), find \( \lim\limits_{x\rightarrow-1}f(x) \).}
{\( 2 \)}{\( 1 \)}{\( -1 \)}{does not exist}{}{}}
\LANGpl{
\Question{
Dany jest wykres funkcji \( f \), Wyznacz \( \lim\limits_{x\rightarrow-1}f(x) \).}
{\( 2 \)}{\( 1 \)}{\( -1 \)}{nie istnieje}{}{}}
\LANGcs{
\Question{
Je dán graf funkce \( f \). Určete \( \lim\limits_{x\rightarrow-1}f(x) \).}
{\( 2 \)}{\( 1 \)}{\( -1 \)}{neexistuje}{}{}}
\LANGsk{
\Question{
Je daný graf funkcie \( f \). Určte \( \lim\limits_{x\rightarrow-1}f(x) \).}
{\( 2 \)}{\( 1 \)}{\( -1 \)}{neexistuje}{}{}}
\LANGes{
\Question{
Dado el gráfico de la función \( f \), halla \( \lim\limits_{x\rightarrow-1}f(x) \).}
{\( 2 \)}{\( 1 \)}{\( -1 \)}{no existe}{}{}}
\End

\Data\ID{1103080001}
\Updated{2020-07-15 03:07:12}
\Level{A}
\SubArea{Limits and continuity}
\IMG{1103080001_0.pdf}
\LANGen{
\Question{
The graph of the function \( f \) is given in the figure. Choose the incorrect statement. Dashed lines represent asymptotes of the function $f$.}
{\( \lim\limits_{x\rightarrow \infty} f(x) = -\infty \)}{\( \lim\limits_{x\rightarrow -2^-} f(x) = -\infty \)}{\( \lim\limits_{x\rightarrow \infty} f(x) = -2 \)}{\( \lim\limits_{x\rightarrow -2} f(x) \) does not exist}{}{}}
\LANGpl{
\Question{
Wykres funkcji  \( f \) przedstawiono. Wybierz fałszywe stwierdzenie. Linie przerywane reprezentują asymptoty funkcji $ f $.}
{\( \lim\limits_{x\rightarrow \infty} f(x) = -\infty \)}{\( \lim\limits_{x\rightarrow -2^-} f(x) = -\infty \)}{\( \lim\limits_{x\rightarrow \infty} f(x) = -2 \)}{\( \lim\limits_{x\rightarrow -2} f(x) \) nie istnieje}{}{}}
\LANGcs{
\Question{
Je dán graf funkce \( f \). Vyberte nepravdivé tvrzení. Čárkované čáry představují asymptoty funkce $f$.}
{\( \lim\limits_{x\rightarrow \infty} f(x) = -\infty \)}{\( \lim\limits_{x\rightarrow -2^-} f(x) = -\infty \)}{\( \lim\limits_{x\rightarrow \infty} f(x) = -2 \)}{\( \lim\limits_{x\rightarrow -2} f(x) \) neexistuje}{}{}}
\LANGsk{
\Question{
Je daný graf funkcie \( f \). Vyberte nesprávne tvrdenie. Prerušované čiary predstavujú asymptoty funkcie $ f $.}
{\( \lim\limits_{x\rightarrow \infty} f(x) = -\infty \)}{\( \lim\limits_{x\rightarrow -2^-} f(x) = -\infty \)}{\( \lim\limits_{x\rightarrow \infty} f(x) = -2 \)}{\( \lim\limits_{x\rightarrow -2} f(x) \) neexistuje}{}{}}
\LANGes{
\Question{
Dado el gráfico de la función \( f \). Elige la proposición falsa. Las líneas discontinuas representan las asíntotas de la función $f$.}
{\( \lim\limits_{x\rightarrow \infty} f(x) = -\infty \)}{\( \lim\limits_{x\rightarrow -2^-} f(x) = -\infty \)}{\( \lim\limits_{x\rightarrow \infty} f(x) = -2 \)}{\( \lim\limits_{x\rightarrow -2} f(x) \) no existe}{}{}}
\End

\Data\ID{1103080002}
\Updated{2020-07-15 03:07:12}
\Level{A}
\SubArea{Limits and continuity}
\IMG{1103080002_0.pdf}
\LANGen{
\Question{
The graph of the function \( f \) is given in the figure. Choose the incorrect statement.}
{\( \lim\limits_{x\rightarrow-1}f(x) \) does not exist}{\( \lim\limits_{x\rightarrow\infty} f(x) = \infty \)}{\( \lim\limits_{x\rightarrow0} f(x) = 0 \)}{\( \lim\limits_{x\rightarrow-\infty} f(x) = 1 \)}{}{}}
\LANGpl{
\Question{
Wykres funkcji  \( f \) przedstawiono. Wybierz fałszywe stwierdzenie.}
{\( \lim\limits_{x\rightarrow-1}f(x) \) nie istnieje}{\( \lim\limits_{x\rightarrow\infty} f(x) = \infty \)}{\( \lim\limits_{x\rightarrow0} f(x) = 0 \)}{\( \lim\limits_{x\rightarrow-\infty} f(x) = 1 \)}{}{}}
\LANGcs{
\Question{
Je dán graf funkce \( f \). Vyberte nepravdivé tvrzení.}
{\( \lim\limits_{x\rightarrow-1}f(x) \) neexistuje}{\( \lim\limits_{x\rightarrow\infty} f(x) = \infty \)}{\( \lim\limits_{x\rightarrow0} f(x) = 0 \)}{\( \lim\limits_{x\rightarrow-\infty} f(x) = 1 \)}{}{}}
\LANGsk{
\Question{
Je daný graf funkcie \( f \). Vyberte nesprávne tvrdenie.}
{\( \lim\limits_{x\rightarrow-1}f(x) \) neexistuje}{\( \lim\limits_{x\rightarrow\infty} f(x) = \infty \)}{\( \lim\limits_{x\rightarrow0} f(x) = 0 \)}{\( \lim\limits_{x\rightarrow-\infty} f(x) = 1 \)}{}{}}
\LANGes{
\Question{
Dado el gráfico de la función \( f \). Elige la proposición falsa.}
{\( \lim\limits_{x\rightarrow-1}f(x) \) no existe}{\( \lim\limits_{x\rightarrow\infty} f(x) = \infty \)}{\( \lim\limits_{x\rightarrow0} f(x) = 0 \)}{\( \lim\limits_{x\rightarrow-\infty} f(x) = 1 \)}{}{}}
\End

\Data\ID{1103080003}
\Updated{2020-07-15 03:07:12}
\Level{A}
\SubArea{Limits and continuity}
\IMG{1103080003.pdf}
\LANGen{
\Question{
The graph of the function \( f \) is given in the figure. Choose the incorrect statement.}
{\( \lim\limits_{x\rightarrow \infty} f(x) = -x \)}{\( \lim\limits_{x\rightarrow 0^+} f(x) = 0 \)}{\( \lim\limits_{x\rightarrow 0^-} f(x) = \infty \)}{\( \lim\limits_{x\rightarrow-\infty} f(x) = \infty \)}{}{}}
\LANGpl{
\Question{
Wykres funkcji  \( f \) przedstawiono. Wybierz fałszywe stwierdzenie.}
{\( \lim\limits_{x\rightarrow \infty} f(x) = -x \)}{\( \lim\limits_{x\rightarrow 0^+} f(x) = 0 \)}{\( \lim\limits_{x\rightarrow 0^-} f(x) = \infty \)}{\( \lim\limits_{x\rightarrow-\infty} f(x) = \infty \)}{}{}}
\LANGcs{
\Question{
Je dán graf funkce \( f \). Vyberte nepravdivé tvrzení.}
{\( \lim\limits_{x\rightarrow \infty} f(x) = -x \)}{\( \lim\limits_{x\rightarrow 0^+} f(x) = 0 \)}{\( \lim\limits_{x\rightarrow 0^-} f(x) = \infty \)}{\( \lim\limits_{x\rightarrow-\infty} f(x) = \infty \)}{}{}}
\LANGsk{
\Question{
Je daný graf funkcie \( f \). Vyberte nesprávne tvrdenie.}
{\( \lim\limits_{x\rightarrow \infty} f(x) = -x \)}{\( \lim\limits_{x\rightarrow 0^+} f(x) = 0 \)}{\( \lim\limits_{x\rightarrow 0^-} f(x) = \infty \)}{\( \lim\limits_{x\rightarrow-\infty} f(x) = \infty \)}{}{}}
\LANGes{
\Question{
Dado el gráfico de la función \( f \). Elige la proposición falsa.}
{\( \lim\limits_{x\rightarrow \infty} f(x) = -x \)}{\( \lim\limits_{x\rightarrow 0^+} f(x) = 0 \)}{\( \lim\limits_{x\rightarrow 0^-} f(x) = \infty \)}{\( \lim\limits_{x\rightarrow-\infty} f(x) = \infty \)}{}{}}
\End

\Data\ID{1103080004}
\Updated{2020-07-15 03:07:12}
\Level{A}
\SubArea{Limits and continuity}
\IMG{1103080004.pdf}
\LANGen{
\Question{
The graph of the function \( f \) is given in the figure. Choose the incorrect statement.}
{\( \lim\limits_{x\rightarrow1^-} f(x) = -1 \)}{\( \lim\limits_{x\rightarrow -1} f(x) \) does not exist}{\( \lim\limits_{x\rightarrow1^+} f(x) = 0 \)}{\( \lim\limits_{x\rightarrow-\infty} f(x) = 1 \)}{}{}}
\LANGpl{
\Question{
Wykres funkcji  \( f \) przedstawiono. Wybierz fałszywe stwierdzenie.}
{\( \lim\limits_{x\rightarrow1^-} f(x) = -1 \)}{\( \lim\limits_{x\rightarrow -1} f(x) \) nie istnieje}{\( \lim\limits_{x\rightarrow1^+} f(x) = 0 \)}{\( \lim\limits_{x\rightarrow-\infty} f(x) = 1 \)}{}{}}
\LANGcs{
\Question{
Je dán graf funkce \( f \). Vyberte nepravdivé tvrzení.}
{\( \lim\limits_{x\rightarrow1^-} f(x) = -1 \)}{\( \lim\limits_{x\rightarrow -1} f(x) \) neexistuje}{\( \lim\limits_{x\rightarrow1^+} f(x) = 0 \)}{\( \lim\limits_{x\rightarrow-\infty} f(x) = 1 \)}{}{}}
\LANGsk{
\Question{
Je daný graf funkcie \( f \). Vyberte nesprávne tvrdenie.}
{\( \lim\limits_{x\rightarrow1^-} f(x) = -1 \)}{\( \lim\limits_{x\rightarrow -1} f(x) \) neexistuje}{\( \lim\limits_{x\rightarrow1^+} f(x) = 0 \)}{\( \lim\limits_{x\rightarrow-\infty} f(x) = 1 \)}{}{}}
\LANGes{
\Question{
Dado el gráfico de la función \( f \). Elige la proposición falsa.}
{\( \lim\limits_{x\rightarrow1^-} f(x) = -1 \)}{\( \lim\limits_{x\rightarrow -1} f(x) \) no existe}{\( \lim\limits_{x\rightarrow1^+} f(x) = 0 \)}{\( \lim\limits_{x\rightarrow-\infty} f(x) = 1 \)}{}{}}
\End

\Data\ID{1103080005}
\Updated{2020-07-15 03:07:21}
\Level{A}
\SubArea{Limits and continuity}
\IMG{1103080005.pdf}
\LANGen{
\Question{
The graph of the function \( f \) is given in the figure bellow. Choose the incorrect statement.}
{\( \lim\limits_{x\to-\infty}f(x)=-1 \)}{\( \lim\limits_{x\to\infty}f(x)=\infty \)}{\( \lim\limits_{x\to1^-}f(x)=2 \)}{\( \lim\limits_{x\to-1^+}f(x)=-\infty \)}{}{}}
\LANGpl{
\Question{
Wykres funkcji  \( f \) przedstawiono poniżej. Wybierz fałszywe stwierdzenie.}
{\( \lim\limits_{x\to-\infty}f(x)=-1 \)}{\( \lim\limits_{x\to\infty}f(x)=\infty \)}{\( \lim\limits_{x\to1^-}f(x)=2 \)}{\( \lim\limits_{x\to-1^+}f(x)=-\infty \)}{}{}}
\LANGcs{
\Question{
Je dán graf funkce \( f \). Vyberte nepravdivé tvrzení.}
{\( \lim\limits_{x\to-\infty}f(x)=-1 \)}{\( \lim\limits_{x\to\infty}f(x)=\infty \)}{\( \lim\limits_{x\to1^-}f(x)=2 \)}{\( \lim\limits_{x\to-1^+}f(x)=-\infty \)}{}{}}
\LANGsk{
\Question{
Je daný graf funkcie \( f \). Vyberte nesprávne tvrdenie.}
{\( \lim\limits_{x\to-\infty}f(x)=-1 \)}{\( \lim\limits_{x\to\infty}f(x)=\infty \)}{\( \lim\limits_{x\to1^-}f(x)=2 \)}{\( \lim\limits_{x\to-1^+}f(x)=-\infty \)}{}{}}
\LANGes{
\Question{
Dado el gráfico de la función \( f \). Elige la proposición falsa.}
{\( \lim\limits_{x\to-\infty}f(x)=-1 \)}{\( \lim\limits_{x\to\infty}f(x)=\infty \)}{\( \lim\limits_{x\to1^-}f(x)=2 \)}{\( \lim\limits_{x\to-1^+}f(x)=-\infty \)}{}{}}
\End

\Data\ID{9000062403}
\Updated{2020-07-15 03:07:37}
\Level{A}
\SubArea{Limits and continuity}
\LANGen{
\Question{
Evaluate the following limit. 
\[ 
 \lim _{x\to -1}\frac{x^{2} - 3x - 4} 
{x^{2} + 6x + 5}
\]}
{\(-\frac{5} 
{4}\)}{\(\frac{4}
{5}\)}{\(\frac{5}
{4}\)}{\(-\frac{4} 
{5}\)}{}{}}
\LANGpl{
\Question{
Oblicz granicę funkcji.
\[ 
 \lim _{x\to -1}\frac{x^{2} - 3x - 4} 
{x^{2} + 6x + 5}
\]}
{\(-\frac{5} 
{4}\)}{\(\frac{4}
{5}\)}{\(\frac{5}
{4}\)}{\(-\frac{4} 
{5}\)}{}{}}
\LANGcs{
\Question{
Určete limitu \(\lim _{x\to -1}\frac{x^{2}-3x-4} 
{x^{2}+6x+5}\).}
{\(-\frac{5} 
{4}\)}{\(\frac{4}
{5}\)}{\(\frac{5}
{4}\)}{\(-\frac{4} 
{5}\)}{}{}}
\LANGsk{
\Question{
Určte limitu \(\lim _{x\to -1}\frac{x^{2}-3x-4} 
{x^{2}+6x+5}\).}
{\(-\frac{5} 
{4}\)}{\(\frac{4}
{5}\)}{\(\frac{5}
{4}\)}{\(-\frac{4} 
{5}\)}{}{}}
\LANGes{
\Question{
Resuelve el siguiente límite:
\[ 
 \lim _{x\to -1}\frac{x^{2} - 3x - 4} 
{x^{2} + 6x + 5}
\]}
{\(-\frac{5} 
{4}\)}{\(\frac{4}
{5}\)}{\(\frac{5}
{4}\)}{\(-\frac{4} 
{5}\)}{}{}}
\End

\Data\ID{9000062404}
\Updated{2020-07-15 03:07:37}
\Level{A}
\SubArea{Limits and continuity}
\LANGen{
\Question{
Evaluate the following limit. 
\[ 
 \lim _{x\to +\infty } \frac{x^{3} - x + 1} 
{1 - x^{2} - x^{3}}
\]}
{\(- 1\)}{\(0.5\)}{\(- 0.5\)}{\(1\)}{}{}}
\LANGpl{
\Question{
Oblicz granicę funkcji. 
\[ 
 \lim _{x\to +\infty } \frac{x^{3} - x + 1} 
{1 - x^{2} - x^{3}}
\]}
{\(- 1\)}{\(0.5\)}{\(- 0.5\)}{\(1\)}{}{}}
\LANGcs{
\Question{
Určete limitu \(\lim _{x\to +\infty } \frac{x^{3}-x+1} 
{1-x^{2}-x^{3}} \).}
{\(- 1\)}{\(0{,}5\)}{\(- 0{,}5\)}{\(1\)}{}{}}
\LANGsk{
\Question{
Určte limitu \(\lim _{x\to +\infty } \frac{x^{3}-x+1} 
{1-x^{2}-x^{3}} \).}
{\(- 1\)}{\(0{,}5\)}{\(- 0{,}5\)}{\(1\)}{}{}}
\LANGes{
\Question{
Resuelve el siguiente límite:
\[ 
 \lim _{x\to +\infty } \frac{x^{3} - x + 1} 
{1 - x^{2} - x^{3}}
\]}
{\(- 1\)}{\(0.5\)}{\(- 0.5\)}{\(1\)}{}{}}
\End

\Data\ID{9000062405}
\Updated{2020-07-15 03:07:37}
\Level{A}
\SubArea{Limits and continuity}
\LANGen{
\Question{
Evaluate the following one-sided limit. 
\[ 
 \lim _{x\to 6^{-}}\frac{3x + 2} 
 {x - 6}
\]}
{\(-\infty \)}{\(1\)}{\(+\infty \)}{\(0\)}{}{}}
\LANGpl{
\Question{
Oblicz jednostronną granicę funkcji.
\[ 
 \lim _{x\to 6^{-}}\frac{3x + 2} 
 {x - 6}
\]}
{\(-\infty \)}{\(1\)}{\(+\infty \)}{\(0\)}{}{}}
\LANGcs{
\Question{
Určete jednostrannou limitu \(\lim _{x\to 6^{-}}\frac{3x+2} 
 {x-6} \).}
{\(-\infty \)}{\(1\)}{\(+\infty \)}{\(0\)}{}{}}
\LANGsk{
\Question{
Určte jednostrannú limitu \(\lim _{x\to 6^{-}}\frac{3x+2} 
 {x-6} \).}
{\(-\infty \)}{\(1\)}{\(+\infty \)}{\(0\)}{}{}}
\LANGes{
\Question{
Resuelve el siguiente límite unilateral:
\[ 
 \lim _{x\to 6^{-}}\frac{3x + 2} 
 {x - 6}
\]}
{\(-\infty \)}{\(1\)}{\(+\infty \)}{\(0\)}{}{}}
\End

\Data\ID{9000141901}
\Updated{2020-07-15 03:07:37}
\Level{A}
\SubArea{Limits and continuity}
\IMG{001419_01-figure0.pdf}
\LANGen{
\Question{
Given the function \(f\),
find \(\lim _{x\to 1}f(x)\).
\[
f(x)=\begin{cases}
 x^3+1 & \text{if } x\neq 1,\\
 3 & \text{if } x = 1
 \end{cases}
\]}
{\(2\)}{\(3\)}{\(1\)}{Does not exist}{}{}}
\LANGpl{
\Question{
Podano funkcję \(f\).
Wyznacz  \(\lim _{x\to 1}f(x)\).
\[
f(x)=\begin{cases}
 x^3+1 & \text{jeśli } x\neq 1,\\
 3 & \text{jeśli } x = 1
 \end{cases}
\]}
{\(2\)}{\(3\)}{\(1\)}{Nie istnieje}{}{}}
\LANGcs{
\Question{
Je dána funkce \(f\) (viz
obrázek). Určete \(\lim _{x\to 1}f(x)\).
\[
f(x)=\begin{cases}
 x^3+1 & \text{pro } x\neq 1,\\
 3 & \text{pro } x = 1
 \end{cases}
\]}
{\(2\)}{\(3\)}{\(1\)}{Limita neexistuje}{}{}}
\LANGsk{
\Question{
Je daná funkcia \(f\) (viď
obrázok). Určte \(\lim _{x\to 1}f(x)\).
\[
f(x)=\begin{cases}
 x^3+1 & \text{pre } x\neq 1,\\
 3 & \text{pre } x = 1
 \end{cases}
\]}
{\(2\)}{\(3\)}{\(1\)}{Limita neexistuje}{}{}}
\LANGes{
\Question{
Dada la función \(f\),
halla \(\lim _{x\to 1}f(x)\).
\[
f(x)=\begin{cases}
 x^3+1 & \text{si } x\neq 1,\\
 3 & \text{si } x = 1
 \end{cases}
\]}
{\(2\)}{\(3\)}{\(1\)}{no existe}{}{}}
\End

\Data\ID{9000141902}
\Updated{2020-07-15 03:07:37}
\Level{A}
\SubArea{Limits and continuity}
\IMG{001419_02-figure0.pdf}
\LANGen{
\Question{
Given the function \(f\),
find \(\lim _{x\to \infty }f(x)\).
\[
f(x)=\begin{cases}
 x^3+1 & \text{if } x\neq 1,\\
 3 & \text{if } x = 1
 \end{cases}
\]}
{\(\infty \)}{\(-\infty \)}{\(4\)}{Does not exist}{}{}}
\LANGpl{
\Question{
Podano funkcję \(f\).
Wyznacz \(\lim _{x\to \infty }f(x)\).
\[
f(x)=\begin{cases}
 x^3+1 & \text{jeśli } x\neq 1,\\
 3 & \text{jeśli } x = 1
 \end{cases}
\]}
{\(\infty \)}{\(-\infty \)}{\(4\)}{Nie istnieje}{}{}}
\LANGcs{
\Question{
Je dána funkce \(f\) (viz
obrázek). Určete \(\lim _{x\to \infty }f(x)\).
\[
f(x)=\begin{cases}
 x^3+1 & \text{pro } x\neq 1,\\
 3 & \text{pro } x = 1
 \end{cases}
\]}
{\(\infty \)}{\(-\infty \)}{\(4\)}{Limita neexistuje}{}{}}
\LANGsk{
\Question{
Je daná funkcia \(f\) (viď
obrázok). Určte \(\lim _{x\to \infty }f(x)\).
\[
f(x)=\begin{cases}
 x^3+1 & \text{pre } x\neq 1,\\
 3 & \text{pre } x = 1
 \end{cases}
\]}
{\(\infty \)}{\(-\infty \)}{\(4\)}{Limita neexistuje}{}{}}
\LANGes{
\Question{
Dada la función \(f\),
halla \(\lim _{x\to \infty }f(x)\).
\[
f(x)=\begin{cases}
 x^3+1 & \text{si } x\neq 1,\\
 3 & \text{si } x = 1
 \end{cases}
\]}
{\(\infty \)}{\(-\infty \)}{\(4\)}{no existe}{}{}}
\End

\Data\ID{9000141903}
\Updated{2020-07-15 03:07:37}
\Level{A}
\SubArea{Limits and continuity}
\IMG{001419_03-figure0.pdf}
\LANGen{
\Question{
Given the function \(g\),
find \(\lim _{x\to 1^{-}}g(x)\).
\[
g(x)=\begin{cases}
 -\frac12(x-1)^2+2 & \text{if } x < 1,\\
 \frac2{x^2}+1 & \text{if } x \geq 1
 \end{cases}
\]}
{\(2\)}{\(3\)}{\(1\)}{Does not exist}{}{}}
\LANGpl{
\Question{
Podano funkcję \(g\),
wyznacz  \(\lim _{x\to 1^{-}}g(x)\).
\[
g(x)=\begin{cases}
 -\frac12(x-1)^2+2 & \text{jeśli } x < 1,\\
 \frac2{x^2}+1 & \text{jeśli } x \geq 1
 \end{cases}
\]}
{\(2\)}{\(3\)}{\(1\)}{Nie istnieje}{}{}}
\LANGcs{
\Question{
Je dána funkce \(g\) (viz
obrázek). Určete \(\lim _{x\to 1^{-}}g(x)\).
\[
g(x)=\begin{cases}
 -\frac12(x-1)^2+2 & \text{pro } x < 1,\\
 \frac2{x^2}+1 & \text{pro } x \geq 1
 \end{cases}
\]}
{\(2\)}{\(3\)}{\(1\)}{Limita neexistuje}{}{}}
\LANGsk{
\Question{
Je daná funkcia \(g\) (viď
obrázok). Určte \(\lim _{x\to 1^{-}}g(x)\).
\[
g(x)=\begin{cases}
 -\frac12(x-1)^2+2 & \text{pre } x < 1,\\
 \frac2{x^2}+1 & \text{pre } x \geq 1
 \end{cases}
\]}
{\(2\)}{\(3\)}{\(1\)}{Limita neexistuje}{}{}}
\LANGes{
\Question{
Dada la función \(g\),
halla \(\lim _{x\to 1^{-}}g(x)\).
\[
g(x)=\begin{cases}
 -\frac12(x-1)^2+2 & \text{si } x < 1,\\
 \frac2{x^2}+1 & \text{si } x \geq 1
 \end{cases}
\]}
{\(2\)}{\(3\)}{\(1\)}{no existe}{}{}}
\End

\Data\ID{9000141904}
\Updated{2020-07-15 03:07:37}
\Level{A}
\SubArea{Limits and continuity}
\IMG{001419_04-figure0.pdf}
\LANGen{
\Question{
Given the function \(g\),
find \(\lim _{x\to 1^{+}}g(x)\).
\[
g(x)=\begin{cases}
 -\frac12(x-1)^2+2 & \text{if } x < 1,\\
 \frac2{x^2}+1 & \text{if } x \geq 1
 \end{cases}
\]}
{\(3\)}{\(2\)}{\(1\)}{Does not exist}{}{}}
\LANGpl{
\Question{
Podano funkcję \(g\),
wyznacz \(\lim _{x\to 1^{+}}g(x)\).
\[
g(x)=\begin{cases}
 -\frac12(x-1)^2+2 & \text{jeśli } x < 1,\\
 \frac2{x^2}+1 & \text{jeśli } x \geq 1
 \end{cases}
\]}
{\(3\)}{\(2\)}{\(1\)}{Nie istnieje.}{}{}}
\LANGcs{
\Question{
Je dána funkce \(g\) (viz
obrázek). Určete \(\lim _{x\to 1^{+}}g(x)\).
\[
g(x)=\begin{cases}
 -\frac12(x-1)^2+2 & \text{pro } x < 1,\\
 \frac2{x^2}+1 & \text{pro } x \geq 1
 \end{cases}
\]}
{\(3\)}{\(2\)}{\(1\)}{Limita neexistuje}{}{}}
\LANGsk{
\Question{
Je daná funkcia \(g\) (viď
obrázok). Určte \(\lim _{x\to 1^{+}}g(x)\).
\[
g(x)=\begin{cases}
 -\frac12(x-1)^2+2 & \text{pre } x < 1,\\
 \frac2{x^2}+1 & \text{pre } x \geq 1
 \end{cases}
\]}
{\(3\)}{\(2\)}{\(1\)}{Limita neexistuje}{}{}}
\LANGes{
\Question{
Dada la función \(g\),
halla \(\lim _{x\to 1^{+}}g(x)\).
\[
g(x)=\begin{cases}
 -\frac12(x-1)^2+2 & \text{si } x < 1,\\
 \frac2{x^2}+1 & \text{si } x \geq 1
 \end{cases}
\]}
{\(3\)}{\(2\)}{\(1\)}{no existe}{}{}}
\End

\Data\ID{9000141905}
\Updated{2020-07-15 03:07:37}
\Level{A}
\SubArea{Limits and continuity}
\IMG{001419_05-figure0.pdf}
\LANGen{
\Question{
Given the function \(g\),
find \(\lim _{x\to 1}g(x)\).
\[
g(x)=\begin{cases}
 -\frac12(x-1)^2+2 & \text{if } x < 1,\\
 \frac2{x^2}+1 & \text{if } x \geq 1
 \end{cases}
\]}
{Does not exist}{\(3\)}{\(2\)}{\(1\)}{}{}}
\LANGpl{
\Question{
Podano funkcję \(g\),
wyznacz \(\lim _{x\to 1}g(x)\).
\[
g(x)=\begin{cases}
 -\frac12(x-1)^2+2 & \text{jeśli } x < 1,\\
 \frac2{x^2}+1 & \text{jeśli } x \geq 1
 \end{cases}
\]}
{Nie istnieje.}{\(3\)}{\(2\)}{\(1\)}{}{}}
\LANGcs{
\Question{
Je dána funkce \(g\) (viz
obrázek). Určete \(\lim _{x\to 1}g(x)\).
\[
g(x)=\begin{cases}
 -\frac12(x-1)^2+2 & \text{pro } x < 1,\\
 \frac2{x^2}+1 & \text{pro } x \geq 1
 \end{cases}
\]}
{Limita neexistuje}{\(3\)}{\(2\)}{\(1\)}{}{}}
\LANGsk{
\Question{
Je daná funkcia \(g\) (viď
obrázok). Určte \(\lim _{x\to 1}g(x)\).
\[
g(x)=\begin{cases}
 -\frac12(x-1)^2+2 & \text{pre } x < 1,\\
 \frac2{x^2}+1 & \text{pre } x \geq 1
 \end{cases}
\]}
{Limita neexistuje}{\(3\)}{\(2\)}{\(1\)}{}{}}
\LANGes{
\Question{
Dada la función \(g\),
halla \(\lim _{x\to 1}g(x)\).
\[
g(x)=\begin{cases}
 -\frac12(x-1)^2+2 & \text{si } x < 1,\\
 \frac2{x^2}+1 & \text{si } x \geq 1
 \end{cases}
\]}
{no existe}{\(3\)}{\(2\)}{\(1\)}{}{}}
\End

\Data\ID{9000141906}
\Updated{2020-07-15 03:07:37}
\Level{A}
\SubArea{Limits and continuity}
\IMG{001419_06-figure0.pdf}
\LANGen{
\Question{
Given the function \(g\),
find \(\lim _{x\to \infty }g(x)\).
\[
g(x)=\begin{cases}
 -\frac12(x-1)^2+2 & \text{if } x < 1,\\
 \frac2{x^2}+1 & \text{if } x \geq 1
 \end{cases}
\]}
{\(1\)}{\(0\)}{\(\infty \)}{\(-\infty \)}{Does not exist}{}}
\LANGpl{
\Question{
Podano funkcję \(g\),
wyznacz \(\lim _{x\to \infty }g(x)\).
\[
g(x)=\begin{cases}
 -\frac12(x-1)^2+2 & \text{jeśli } x < 1,\\
 \frac2{x^2}+1 & \text{jeśli } x \geq 1
 \end{cases}
\]}
{\(1\)}{\(0\)}{\(\infty \)}{\(-\infty \)}{Nie istnieje}{}}
\LANGcs{
\Question{
Je dána funkce \(g\) (viz
obrázek). Určete \(\lim _{x\to \infty }g(x)\).
\[
g(x)=\begin{cases}
 -\frac12(x-1)^2+2 & \text{pro } x < 1,\\
 \frac2{x^2}+1 & \text{pro } x \geq 1
 \end{cases}
\]}
{\(1\)}{\(0\)}{\(\infty \)}{\(-\infty \)}{Limita neexistuje}{}}
\LANGsk{
\Question{
Je daná funkcia \(g\) (viď
obrázok). Určte \(\lim _{x\to \infty }g(x)\).
\[
g(x)=\begin{cases}
 -\frac12(x-1)^2+2 & \text{pre } x < 1,\\
 \frac2{x^2}+1 & \text{pre } x \geq 1
 \end{cases}
\]}
{\(1\)}{\(0\)}{\(\infty \)}{\(-\infty \)}{Limita neexistuje}{}}
\LANGes{
\Question{
Dada la función \(g\),
halla \(\lim _{x\to \infty }g(x)\).
\[
g(x)=\begin{cases}
 -\frac12(x-1)^2+2 & \text{si } x < 1,\\
 \frac2{x^2}+1 & \text{si } x \geq 1
 \end{cases}
\]}
{\(1\)}{\(0\)}{\(\infty \)}{\(-\infty \)}{no existe}{}}
\End

\Data\ID{9000141907}
\Updated{2020-07-15 03:07:37}
\Level{A}
\SubArea{Limits and continuity}
\IMG{001419_07-figure0.pdf}
\LANGen{
\Question{
Given the function \(h\),
find \(\lim _{x\to 1^{-}}h(x)\).
\[
h(x)=\begin{cases}
 -\frac1{x-1} & \text{if } x< 1,\\
 -(x-1)^2+2 & \text{if } x\geq 1
 \end{cases}
\]}
{\(\infty \)}{\(1\)}{\(2\)}{\(-\infty \)}{Does not exist}{}}
\LANGpl{
\Question{
Podano funkcję \(h\),
wyznacz \(\lim _{x\to 1^{-}}h(x)\).
\[
h(x)=\begin{cases}
 -\frac1{x-1} & \text{jeśli } x< 1,\\
 -(x-1)^2+2 & \text{jeśli } x\geq 1
 \end{cases}
\]}
{\(\infty \)}{\(1\)}{\(2\)}{\(-\infty \)}{Nie istnieje.}{}}
\LANGcs{
\Question{
Je dána funkce \(h\) (viz
obrázek). Určete \(\lim _{x\to 1^{-}}h(x)\).
\[
h(x)=\begin{cases}
 -\frac1{x-1} & \text{pro } x< 1,\\
 -(x-1)^2+2 & \text{pro } x\geq 1
 \end{cases}
\]}
{\(\infty \)}{\(1\)}{\(2\)}{\(-\infty \)}{Limita neexistuje}{}}
\LANGsk{
\Question{
Je daná funkcia \(h\) (viď
obrázok). Určte \(\lim _{x\to 1^{-}}h(x)\).
\[
h(x)=\begin{cases}
 -\frac1{x-1} & \text{pre } x< 1,\\
 -(x-1)^2+2 & \text{pre } x\geq 1
 \end{cases}
\]}
{\(\infty \)}{\(1\)}{\(2\)}{\(-\infty \)}{Limita neexistuje}{}}
\LANGes{
\Question{
Dada la función \(h\),
halla \(\lim _{x\to 1^{-}}h(x)\).
\[
h(x)=\begin{cases}
 -\frac1{x-1} & \text{si } x< 1,\\
 -(x-1)^2+2 & \text{si } x\geq 1
 \end{cases}
\]}
{\(\infty \)}{\(1\)}{\(2\)}{\(-\infty \)}{no existe}{}}
\End

\Data\ID{9000141908}
\Updated{2020-07-15 03:07:37}
\Level{A}
\SubArea{Limits and continuity}
\IMG{001419_08-figure0.pdf}
\LANGen{
\Question{
Given the function \(h\),
find \(\lim _{x\to 1^{+}}h(x)\).
\[
h(x)=\begin{cases}
 -\frac1{x-1} & \text{if } x< 1,\\
 -(x-1)^2+2 & \text{if } x\geq 1
 \end{cases}
\]}
{\(2\)}{\(1\)}{\(0\)}{\(\infty \)}{Does not exist}{}}
\LANGpl{
\Question{
Podano funkcję \(h\),
wyznacz \(\lim _{x\to 1^{+}}h(x)\).
\[
h(x)=\begin{cases}
 -\frac1{x-1} & \text{jeśli } x< 1,\\
 -(x-1)^2+2 & \text{jeśli } x\geq 1
 \end{cases}
\]}
{\(2\)}{\(1\)}{\(0\)}{\(\infty \)}{Nie istnieje.}{}}
\LANGcs{
\Question{
Je dána funkce \(h\) (viz
obrázek). Určete \(\lim _{x\to 1^{+}}h(x)\).
\[
h(x)=\begin{cases}
 -\frac1{x-1} & \text{pro } x< 1,\\
 -(x-1)^2+2 & \text{pro } x\geq 1
 \end{cases}
\]}
{\(2\)}{\(1\)}{\(0\)}{\(\infty \)}{Limita neexistuje}{}}
\LANGsk{
\Question{
Je daná funkcia \(h\) (viď
obrázok). Určte \(\lim _{x\to 1^{+}}h(x)\).
\[
h(x)=\begin{cases}
 -\frac1{x-1} & \text{pre } x< 1,\\
 -(x-1)^2+2 & \text{pre } x\geq 1
 \end{cases}
\]}
{\(2\)}{\(1\)}{\(0\)}{\(\infty \)}{Limita neexistuje}{}}
\LANGes{
\Question{
Dada la función \(h\),
halla \(\lim _{x\to 1^{+}}h(x)\).
\[
h(x)=\begin{cases}
 -\frac1{x-1} & \text{si } x< 1,\\
 -(x-1)^2+2 & \text{si } x\geq 1
 \end{cases}
\]}
{\(2\)}{\(1\)}{\(0\)}{\(\infty \)}{no existe}{}}
\End

\Data\ID{9000141909}
\Updated{2020-07-15 03:07:37}
\Level{A}
\SubArea{Limits and continuity}
\IMG{001419_09-figure0.pdf}
\LANGen{
\Question{
Given the function \(h\),
find \(\lim _{x\to \infty }h(x)\).
\[
h(x)=\begin{cases}
 -\frac1{x-1} & \text{if } x< 1,\\
 -(x-1)^2+2 & \text{if } x\geq 1
 \end{cases}
\]}
{\(-\infty \)}{\(1\)}{\(0\)}{\(\infty \)}{Does not exist}{}}
\LANGpl{
\Question{
Podano funkcję \(h\),
wyznacz \(\lim _{x\to \infty }h(x)\).
\[
h(x)=\begin{cases}
 -\frac1{x-1} & \text{jeśli } x< 1,\\
 -(x-1)^2+2 & \text{jeśli } x\geq 1
 \end{cases}
\]}
{\(-\infty \)}{\(1\)}{\(0\)}{\(\infty \)}{Nie istnieje.}{}}
\LANGcs{
\Question{
Je dána funkce \(h\) (viz
obrázek). Určete \(\lim _{x\to \infty }h(x)\).
\[
h(x)=\begin{cases}
 -\frac1{x-1} & \text{pro } x< 1,\\
 -(x-1)^2+2 & \text{pro } x\geq 1
 \end{cases}
\]}
{\(-\infty \)}{\(1\)}{\(0\)}{\(\infty \)}{Limita neexistuje}{}}
\LANGsk{
\Question{
Je daná funkcia \(h\) (viď
obrázok). Určte \(\lim _{x\to \infty }h(x)\).
\[
h(x)=\begin{cases}
 -\frac1{x-1} & \text{pre } x< 1,\\
 -(x-1)^2+2 & \text{pre } x\geq 1
 \end{cases}
\]}
{\(-\infty \)}{\(1\)}{\(0\)}{\(\infty \)}{Limita neexistuje}{}}
\LANGes{
\Question{
Dada la función \(h\),
halla \(\lim _{x\to \infty }h(x)\).
\[
h(x)=\begin{cases}
 -\frac1{x-1} & \text{si } x< 1,\\
 -(x-1)^2+2 & \text{si } x\geq 1
 \end{cases}
\]}
{\(-\infty \)}{\(1\)}{\(0\)}{\(\infty \)}{no existe}{}}
\End

\Data\ID{9000141910}
\Updated{2020-07-15 03:07:37}
\Level{A}
\SubArea{Limits and continuity}
\IMG{001419_10-figure0.pdf}
\LANGen{
\Question{
Given the function \(h\),
find \(\lim _{x\to -\infty }h(x)\).
\[
h(x)=\begin{cases}
 -\frac1{x-1} & \text{if } x< 1,\\
 -(x-1)^2+2 & \text{if } x\geq 1
 \end{cases}
\]}
{\(0\)}{\(2\)}{\(\infty \)}{\(-\infty \)}{Does not exist}{}}
\LANGpl{
\Question{
Podano funkcję \(h\),
wyznacz \(\lim _{x\to -\infty }h(x)\).
\[
h(x)=\begin{cases}
 -\frac1{x-1} & \text{jeśli } x< 1,\\
 -(x-1)^2+2 & \text{jeśli } x\geq 1
 \end{cases}
\]}
{\(0\)}{\(2\)}{\(\infty \)}{\(-\infty \)}{Nie istnieje.}{}}
\LANGcs{
\Question{
Je dána funkce \(h\) (viz
obrázek). Určete \(\lim _{x\to -\infty }h(x)\).
\[
h(x)=\begin{cases}
 -\frac1{x-1} & \text{pro } x< 1,\\
 -(x-1)^2+2 & \text{pro } x\geq 1
 \end{cases}
\]}
{\(0\)}{\(2\)}{\(\infty \)}{\(-\infty \)}{Limita neexistuje}{}}
\LANGsk{
\Question{
Je daná funkcia \(h\) (viď
obrázok). Určte \(\lim _{x\to -\infty }h(x)\).
\[
h(x)=\begin{cases}
 -\frac1{x-1} & \text{pre } x< 1,\\
 -(x-1)^2+2 & \text{pre } x\geq 1
 \end{cases}
\]}
{\(0\)}{\(2\)}{\(\infty \)}{\(-\infty \)}{Limita neexistuje}{}}
\LANGes{
\Question{
Dada la función \(h\),
halla \(\lim _{x\to -\infty }h(x)\).
\[
h(x)=\begin{cases}
 -\frac1{x-1} & \text{si } x< 1,\\
 -(x-1)^2+2 & \text{si } x\geq 1
 \end{cases}
\]}
{\(0\)}{\(2\)}{\(\infty \)}{\(-\infty \)}{no existe}{}}
\End

\Data\ID{1003021101}
\Updated{2020-07-15 03:07:12}
\Level{B}
\SubArea{Limits and continuity}
\LANGen{
\Question{
Evaluate the following limit.
\[ \lim_{x\to -2}\frac{2-\sqrt{x+6}}{x-2} \]}
{\( 0 \)}{\( \frac14 \)}{\( -\frac14 \)}{\( 1 \)}{}{}}
\LANGpl{
\Question{
Oblicz granicę funkcji.
\[ \lim_{x\to -2}\frac{2-\sqrt{x+6}}{x-2} \]}
{\( 0 \)}{\( \frac14 \)}{\( -\frac14 \)}{\( 1 \)}{}{}}
\LANGcs{
\Question{
Vypočítejte následující limitu.
\[ \lim_{x\to-2}\frac{2-\sqrt{x+6}}{x-2} \]}
{\( 0 \)}{\( \frac14 \)}{\( -\frac14 \)}{\( 1 \)}{}{}}
\LANGsk{
\Question{
Vypočítajte následujúcu limitu
\[ \lim_{x\to-2}\frac{2-\sqrt{x+6}}{x-2} \]}
{\( 0 \)}{\( \frac14 \)}{\( -\frac14 \)}{\( 1 \)}{}{}}
\LANGes{
\Question{
Resuelve el siguiente límite:
\[ \lim_{x\to -2}\frac{2-\sqrt{x+6}}{x-2} \]}
{\( 0 \)}{\( \frac14 \)}{\( -\frac14 \)}{\( 1 \)}{}{}}
\End

\Data\ID{1003021102}
\Updated{2020-07-15 03:07:12}
\Level{B}
\SubArea{Limits and continuity}
\LANGen{
\Question{
Evaluate the following limit.
\[ \lim_{x\to\frac{\pi}4}\frac{1+\sin2x}{1-\cos4x} \]}
{\( 1 \)}{\( 0 \)}{\( 2 \)}{\( \frac12 \)}{}{}}
\LANGpl{
\Question{
Oblicz granicę funkcji.
\[ \lim_{x\to\frac{\pi}4}\frac{1+\sin2x}{1-\cos4x} \]}
{\( 1 \)}{\( 0 \)}{\( 2 \)}{\( \frac12 \)}{}{}}
\LANGcs{
\Question{
Vypočítejte následující limitu.
\[ \lim_{x\to\frac{\pi}4}\frac{1+\sin2x}{1-\cos4x} \]}
{\( 1 \)}{\( 0 \)}{\( 2 \)}{\( \frac12 \)}{}{}}
\LANGsk{
\Question{
Vypočítajte následujúcu limitu.
\[ \lim_{x\to\frac{\pi}4}\frac{1+\sin2x}{1-\cos4x} \]}
{\( 1 \)}{\( 0 \)}{\( 2 \)}{\( \frac12 \)}{}{}}
\LANGes{
\Question{
Resuelve el siguiente límite:
\[ \lim_{x\to\frac{\pi}4}\frac{1+\sin2x}{1-\cos4x} \]}
{\( 1 \)}{\( 0 \)}{\( 2 \)}{\( \frac12 \)}{}{}}
\End

\Data\ID{1003021107}
\Updated{2020-07-15 03:07:12}
\Level{B}
\SubArea{Limits and continuity}
\LANGen{
\Question{
Evaluate the following limit.
\[ \lim_{x\to\pi}\frac{\mathrm{tg}\,x}{\sin2x} \]}
{\( \frac12 \)}{\( -\frac12 \)}{\( 0 \)}{\( -2 \)}{}{}}
\LANGpl{
\Question{
Oblicz granicę funkcji.
\[ \lim_{x\to\pi}\frac{\mathrm{tg}\,x}{\sin2x} \]}
{\( \frac12 \)}{\( -\frac12 \)}{\( 0 \)}{\( -2 \)}{}{}}
\LANGcs{
\Question{
Vypočítejte následující limitu.
\[ \lim_{x\to\pi}\frac{\mathrm{tg}\,x}{\sin2x} \]}
{\( \frac12 \)}{\( -\frac12 \)}{\( 0 \)}{\( -2 \)}{}{}}
\LANGsk{
\Question{
Vypočítajte následujúcu limitu.
\[ \lim_{x\to\pi}\frac{\mathrm{tg}\,x}{\sin2x} \]}
{\( \frac12 \)}{\( -\frac12 \)}{\( 0 \)}{\( -2 \)}{}{}}
\LANGes{
\Question{
Resuelve el siguiente límite:
\[ \lim_{x\to\pi}\frac{\mathrm{tg}\,x}{\sin2x} \]}
{\( \frac12 \)}{\( -\frac12 \)}{\( 0 \)}{\( -2 \)}{}{}}
\End

\Data\ID{1003021108}
\Updated{2020-07-15 03:07:12}
\Level{B}
\SubArea{Limits and continuity}
\LANGen{
\Question{
Evaluate the following limit.
\[ \lim_{x\to\frac{\pi}4}\frac{\cos x-\sin x}{1-\mathrm{cotg}\,x} \]}
{\( -\frac{\sqrt2}{2} \)}{\( \frac{\sqrt2}{2} \)}{\( 0 \)}{\( -\sqrt2 \)}{}{}}
\LANGpl{
\Question{
Oblicz granicę funkcji.
\[ \lim_{x\to\frac{\pi}4}\frac{\cos x-\sin x}{1-\mathrm{cotg}\,x} \]}
{\( -\frac{\sqrt2}{2} \)}{\( \frac{\sqrt2}{2} \)}{\( 0 \)}{\( -\sqrt2 \)}{}{}}
\LANGcs{
\Question{
Vypočítejte následující limitu.
\[ \lim_{x\to\frac{\pi}4}\frac{\cos x-\sin x}{1-\mathrm{cotg}\,x} \]}
{\( -\frac{\sqrt2}{2} \)}{\( \frac{\sqrt2}{2} \)}{\( 0 \)}{\( -\sqrt2 \)}{}{}}
\LANGsk{
\Question{
Vypočítajte následujúcu limitu.
\[ \lim_{x\to\frac{\pi}4}\frac{\cos x-\sin x}{1-\mathrm{cotg}\,x} \]}
{\( -\frac{\sqrt2}{2} \)}{\( \frac{\sqrt2}{2} \)}{\( 0 \)}{\( -\sqrt2 \)}{}{}}
\LANGes{
\Question{
Resuelve el siguiente límite:
\[ \lim_{x\to\frac{\pi}4}\frac{\cos x-\sin x}{1-\mathrm{cotg}\,x} \]}
{\( -\frac{\sqrt2}{2} \)}{\( \frac{\sqrt2}{2} \)}{\( 0 \)}{\( -\sqrt2 \)}{}{}}
\End

\Data\ID{1003021109}
\Updated{2020-07-15 03:07:12}
\Level{B}
\SubArea{Limits and continuity}
\LANGen{
\Question{
Evaluate the following limit.
\[ \lim\limits_{x\rightarrow0}\frac{\sin2x}{x \cos x} \]}
{\( 2 \)}{\( 1 \)}{\( 0 \)}{\( \frac12 \)}{}{}}
\LANGpl{
\Question{
Oblicz granicę funkcji.
\[ \lim\limits_{x\rightarrow0}\frac{\sin2x}{x \cos x} \]}
{\( 2 \)}{\( 1 \)}{\( 0 \)}{\( \frac12 \)}{}{}}
\LANGcs{
\Question{
Vypočítejte následující limitu.
\[ \lim\limits_{x\rightarrow0}\frac{\sin2x}{x \cos x} \]}
{\( 2 \)}{\( 1 \)}{\( 0 \)}{\( \frac12 \)}{}{}}
\LANGsk{
\Question{
Vypočítajte následujúcu limitu.
\[ \lim\limits_{x\rightarrow0}\frac{\sin2x}{x \cos x} \]}
{\( 2 \)}{\( 1 \)}{\( 0 \)}{\( \frac12 \)}{}{}}
\LANGes{
\Question{
Resuelve el siguiente límite:
\[ \lim\limits_{x\rightarrow0}\frac{\sin2x}{x \cos x} \]}
{\( 2 \)}{\( 1 \)}{\( 0 \)}{\( \frac12 \)}{}{}}
\End

\Data\ID{1003021110}
\Updated{2020-07-15 03:07:12}
\Level{B}
\SubArea{Limits and continuity}
\LANGen{
\Question{
Evaluate the following limit.
\[ \lim\limits_{x\rightarrow6}\frac{x-6}{\sqrt{x+3}-3} \]}
{\( 6 \)}{\( 0 \)}{\( 9 \)}{\( 12 \)}{}{}}
\LANGpl{
\Question{
Oblicz granicę funkcji.
\[ \lim\limits_{x\rightarrow6}\frac{x-6}{\sqrt{x+3}-3} \]}
{\( 6 \)}{\( 0 \)}{\( 9 \)}{\( 12 \)}{}{}}
\LANGcs{
\Question{
Vypočítejte následující limitu.
\[ \lim\limits_{x\rightarrow6}\frac{x-6}{\sqrt{x+3}-3} \]}
{\( 6 \)}{\( 0 \)}{\( 9 \)}{\( 12 \)}{}{}}
\LANGsk{
\Question{
Vypočítajte nasledujúcu limitu.
\[ \lim\limits_{x\rightarrow6}\frac{x-6}{\sqrt{x+3}-3} \]}
{\( 6 \)}{\( 0 \)}{\( 9 \)}{\( 12 \)}{}{}}
\LANGes{
\Question{
Resuelve el siguiente límite:
\[ \lim\limits_{x\rightarrow6}\frac{x-6}{\sqrt{x+3}-3} \]}
{\( 6 \)}{\( 0 \)}{\( 9 \)}{\( 12 \)}{}{}}
\End

\Data\ID{1003021111}
\Updated{2021-01-31 10:01:23}
\Level{B}
\SubArea{Limits and continuity}
\LANGen{
\Question{
Evaluate the following limit.
\[ \lim_{x\to 3}\frac{9-x^2}{2-\sqrt{7-x}} \]}
{\( -24 \)}{\( 0 \)}{\( 24 \)}{\( 36 \)}{}{}}
\LANGpl{
\Question{
Oblicz granicę funkcji.
\[ \lim_{x\to 3}\frac{9-x^2}{2-\sqrt{7-x}} \]}
{\( -24 \)}{\( 0 \)}{\( 24 \)}{\( 36 \)}{}{}}
\LANGcs{
\Question{
Vypočítejte následující limitu.
\[ \lim_{x\to 3}\frac{9-x^2}{2-\sqrt{7-x}} \]}
{\( -24 \)}{\( 0 \)}{\( 24 \)}{\( 36 \)}{}{}}
\LANGsk{
\Question{
Vypočítajte nasledujúcu limitu.
\[ \lim_{x\to 3}\frac{9-x^2}{2-\sqrt{7-x}} \]}
{\( -24 \)}{\( 0 \)}{\( 24 \)}{\( 36 \)}{}{}}
\LANGes{
\Question{
Resuelve el siguiente límite:
\[ \lim_{x\to 3}\frac{9-x^2}{2-\sqrt{7-x}} \]}
{\( -24 \)}{\( 0 \)}{\( 24 \)}{\( 36 \)}{}{}}
\End

\Data\ID{1003021112}
\Updated{2020-07-15 03:07:12}
\Level{B}
\SubArea{Limits and continuity}
\LANGen{
\Question{
Evaluate the following limit.
\[ \lim\limits_{x\rightarrow0}\frac{\sqrt{1+x^2}-1}x \]}
{\( 0 \)}{\( 1 \)}{\( \frac12 \)}{\( 2 \)}{}{}}
\LANGpl{
\Question{
Oblicz granicę funkcji.
\[ \lim\limits_{x\rightarrow0}\frac{\sqrt{1+x^2}-1}x \]}
{\( 0 \)}{\( 1 \)}{\( \frac12 \)}{\( 2 \)}{}{}}
\LANGcs{
\Question{
Vypočítejte následující limitu.
\[ \lim\limits_{x\rightarrow0}\frac{\sqrt{1+x^2}-1}x \]}
{\( 0 \)}{\( 1 \)}{\( \frac12 \)}{\( 2 \)}{}{}}
\LANGsk{
\Question{
Vypočítajte nasledujúcu limitu.
\[ \lim\limits_{x\rightarrow0}\frac{\sqrt{1+x^2}-1}x \]}
{\( 0 \)}{\( 1 \)}{\( \frac12 \)}{\( 2 \)}{}{}}
\LANGes{
\Question{
Resuelve el siguiente límite:
\[ \lim\limits_{x\rightarrow0}\frac{\sqrt{1+x^2}-1}x \]}
{\( 0 \)}{\( 1 \)}{\( \frac12 \)}{\( 2 \)}{}{}}
\End

\Data\ID{1003021113}
\Updated{2020-07-15 03:07:12}
\Level{B}
\SubArea{Limits and continuity}
\LANGen{
\Question{
Evaluate the following limit.
\[ \lim_{x\to0}\frac{\sin \frac x2}x \]}
{\( \frac12 \)}{\( 2 \)}{\( 1 \)}{\( 0 \)}{}{}}
\LANGpl{
\Question{
Oblicz granicę funkcji.
\[ \lim_{x\to0}\frac{\sin \frac x2}x \]}
{\( \frac12 \)}{\( 2 \)}{\( 1 \)}{\( 0 \)}{}{}}
\LANGcs{
\Question{
Vypočítejte následující limitu.
\[ \lim_{x\to0}\frac{\sin \frac x2}x \]}
{\( \frac12 \)}{\( 2 \)}{\( 1 \)}{\( 0 \)}{}{}}
\LANGsk{
\Question{
Vypočítajte nasledujúcu limitu.
\[ \lim_{x\to0}\frac{\sin \frac x2}x \]}
{\( \frac12 \)}{\( 2 \)}{\( 1 \)}{\( 0 \)}{}{}}
\LANGes{
\Question{
Resuelve el siguiente límite:
\[ \lim_{x\to0}\frac{\sin \frac x2}x \]}
{\( \frac12 \)}{\( 2 \)}{\( 1 \)}{\( 0 \)}{}{}}
\End

\Data\ID{1003085501}
\Updated{2020-07-15 03:07:12}
\Level{B}
\SubArea{Limits and continuity}
\LANGen{
\Question{
Decide which of the following functions are continuous at \( x = 1 \).
\[\begin{aligned} f_1(x)&=\frac{x^2+1}{x-1} \\
f_2(x)&=\sqrt{x-1} \\
f_3(x)&=\log x \\
f_4(x)&=\mathrm{tg}(x-1)
\end{aligned}\]
The only such functions are:}
{\( f_3 \), \( f_4 \)}{\( f_2 \), \( f_3 \), \( f_4 \)}{\( f_2 \), \( f_3 \)}{\( f_3 \)}{}{}}
\LANGpl{
\Question{
Które z podanych funkcji to  funkcje ciągłe  w punkcie \( x = 1 \).
\[\begin{aligned} f_1\colon y&=\frac{x^2+1}{x-1} \\
f_2\colon y&=\sqrt{x-1} \\
f_3\colon y&=\log x \\
f_4\colon y&=\mathrm{tg}(x-1)
\end{aligned}\]
Tylko funkcje:}
{\( f_3 \), \( f_4 \)}{\( f_2 \), \( f_3 \), \( f_4 \)}{\( f_2 \), \( f_3 \)}{\( f_3 \)}{}{}}
\LANGcs{
\Question{
Rozhodněte, které z následujících funkcí jsou spojité v bodě  \( x = 1 \).
\[\begin{aligned} f_1\colon y&=\frac{x^2+1}{x-1} \\
f_2\colon y&=\sqrt{x-1} \\
f_3\colon y&=\log x\\
f_4\colon y&=\mathrm{tg}(x-1)
\end{aligned}\]
Jedinými takovými funkcemi jsou:}
{\( f_3 \), \( f_4 \)}{\( f_2 \), \( f_3 \), \( f_4 \)}{\( f_2 \), \( f_3 \)}{\( f_3 \)}{}{}}
\LANGsk{
\Question{
Rozhodnite, ktoré z nasledujúcich funkcií sú spojité v bode  \( x = 1 \).
\[\begin{aligned} f_1\colon y&=\frac{x^2+1}{x-1} \\
f_2\colon y&=\sqrt{x-1} \\
f_3\colon y&=\log x\\
f_4\colon y&=\mathrm{tg}(x-1)
\end{aligned}\]
Jedinými takými funkciami sú:}
{\( f_3 \), \( f_4 \)}{\( f_2 \), \( f_3 \), \( f_4 \)}{\( f_2 \), \( f_3 \)}{\( f_3 \)}{}{}}
\LANGes{
\Question{
Decide cuáles de las siguientes funciones son continuas en \( x = 1 \).
\[\begin{aligned} f_1(x)&=\frac{x^2+1}{x-1} \\
f_2(x)&=\sqrt{x-1} \\
f_3(x)&=\log x \\
f_4(x)&=\mathrm{tg}(x-1)
\end{aligned}\]
Las únicas funciones así son:}
{\( f_3 \), \( f_4 \)}{\( f_2 \), \( f_3 \), \( f_4 \)}{\( f_2 \), \( f_3 \)}{\( f_3 \)}{}{}}
\End

\Data\ID{1003085502}
\Updated{2020-07-15 03:07:12}
\Level{B}
\SubArea{Limits and continuity}
\LANGen{
\Question{
Determine the value of \( a \) (\(a\in\mathbb{R}\)) for which the function \( f(x)=\frac{ax}{x^2+a} \) is not continuous at \( x = -1 \).}
{\( -1 \)}{\( 1 \)}{\( 0 \)}{\( 1\), \( -1 \)}{}{}}
\LANGpl{
\Question{
Dla jakich wartości parametru  \( a \) (\(a\in\mathbb{R}\)) funkcja \( f\colon y=\frac{ax}{x^2+a} \) nie jest  ciągła w punkcie \( x = -1 \).}
{\( -1 \)}{\( 1 \)}{\( 0 \)}{\( 1\), \( -1 \)}{}{}}
\LANGcs{
\Question{
Určete, pro které hodnoty parametru \( a \) (\(a\in\mathbb{R}\)) není funkce \( f\colon y=\frac{ax}{x^2+a} \)  spojitá v bodě \( x = -1 \).}
{\( -1 \)}{\( 1 \)}{\( 0 \)}{\( 1\), \( -1 \)}{}{}}
\LANGsk{
\Question{
Určte, pre ktoré hodnoty parametru \( a \) (\(a\in\mathbb{R}\)) nie je funkcia \( f\colon y=\frac{ax}{x^2+a} \)  spojitá v bode \( x = -1 \).}
{\( -1 \)}{\( 1 \)}{\( 0 \)}{\( 1\), \( -1 \)}{}{}}
\LANGes{
\Question{
Determina para qué valor de \( a \) (\(a\in\mathbb{R}\)) la función \( f(x)=\frac{ax}{x^2+a} \) no es continua en \( x = -1 \).}
{\( -1 \)}{\( 1 \)}{\( 0 \)}{\( 1\), \( -1 \)}{}{}}
\End

\Data\ID{1003085503}
\Updated{2020-07-15 03:07:12}
\Level{B}
\SubArea{Limits and continuity}
\LANGen{
\Question{
Determine the value of \( a \) (\(a\in\mathbb{R}\)) for which the function \( f(x)=\frac{x+a}{ax^2-3} \) is not continuous at \( x = 1 \).}
{\( 3 \)}{\( -3 \)}{\( \sqrt3 \)}{\( -\sqrt3 \)}{}{}}
\LANGpl{
\Question{
Dla jakich wartości parametru \( a \) (\(a\in\mathbb{R}\)) funkcja \( f\colon y=\frac{x+a}{ax^2-3} \) nie jest ciągła w punkcie \( x = 1 \).}
{\( 3 \)}{\( -3 \)}{\( \sqrt3 \)}{\( -\sqrt3 \)}{}{}}
\LANGcs{
\Question{
Určete, pro které hodnoty parametru \( a \) (\(a\in\mathbb{R}\)) není funkce \( f\colon y=\frac{x+a}{ax^2-3} \) spojitá v bodě \( x = 1 \).}
{\( 3 \)}{\( -3 \)}{\( \sqrt3 \)}{\( -\sqrt3 \)}{}{}}
\LANGsk{
\Question{
Určte, pre ktoré hodnoty parametra \( a \) (\(a\in\mathbb{R}\)) nie je funkcia \( f\colon y=\frac{x+a}{ax^2-3} \) spojitá v bode \( x = 1 \).}
{\( 3 \)}{\( -3 \)}{\( \sqrt3 \)}{\( -\sqrt3 \)}{}{}}
\LANGes{
\Question{
Determina para qué valor de \( a \) (\(a\in\mathbb{R}\)) la función \( f(x)=\frac{x+a}{ax^2-3} \) no es continua en \( x = 1 \).}
{\( 3 \)}{\( -3 \)}{\( \sqrt3 \)}{\( -\sqrt3 \)}{}{}}
\End

\Data\ID{1003085504}
\Updated{2020-07-15 03:07:12}
\Level{B}
\SubArea{Limits and continuity}
\LANGen{
\Question{
Determine the value of \( a \) (\(a\in\mathbb{R}\)) for which the function
\[ f(x)=\left\{ \begin{array}{lc} x^2+x & \text{if } x\leq -2 \\ ax+3 & \text{if }x > -2 \end{array}\right. \]
is continuous at \( x = -2 \).}
{\( \frac12 \)}{\( -\frac12 \)}{\( \frac32 \)}{\( -\frac32 \)}{}{}}
\LANGpl{
\Question{
Dla jakich wartości parametru \( a \) (\(a\in\mathbb{R}\)) funkcja
\[ f\colon y=\left\{ \begin{array}{lc} x^2+x & \text{jeśli } x\leq -2 \\ ax+3 & \text{jeśli }x > -2 \end{array}\right. \]
jest ciągła w punkcie \( x = -2 \).}
{\( \frac12 \)}{\( -\frac12 \)}{\( \frac32 \)}{\( -\frac32 \)}{}{}}
\LANGcs{
\Question{
Určete, pro které hodnoty parametru \( a \) (\(a\in\mathbb{R}\)) je funkce
\[ f\colon y=\left\{ \begin{array}{lc} x^2+x & \text{pro } x\leq -2 \\ ax+3 & \text{pro }x > -2 \end{array}\right. \]
spojitá v bodě \( x = -2 \).}
{\( \frac12 \)}{\( -\frac12 \)}{\( \frac32 \)}{\( -\frac32 \)}{}{}}
\LANGsk{
\Question{
Určte, pre ktoré hodnoty parametra \( a \) (\(a\in\mathbb{R}\)) je funkcia
\[ f\colon y=\left\{ \begin{array}{lc} x^2+x & \text{pre } x\leq -2 \\ ax+3 & \text{pre }x > -2 \end{array}\right. \]
spojitá v bode \( x = -2 \).}
{\( \frac12 \)}{\( -\frac12 \)}{\( \frac32 \)}{\( -\frac32 \)}{}{}}
\LANGes{
\Question{
Determina el valor de \( a \) (\(a\in\mathbb{R}\)) para el cual la función
\[ f(x)=\left\{ \begin{array}{lc} x^2+x & \text{if } x\leq -2 \\ ax+3 & \text{if }x > -2 \end{array}\right. \]
es continua en \( x = -2 \).}
{\( \frac12 \)}{\( -\frac12 \)}{\( \frac32 \)}{\( -\frac32 \)}{}{}}
\End

\Data\ID{1003085505}
\Updated{2020-07-15 03:07:12}
\Level{B}
\SubArea{Limits and continuity}
\LANGen{
\Question{
Decide which of the following functions have points of discontinuity.
\[ \begin{aligned}
f_1(x)&=\frac{x-1}{x^2+1} \\
f_2(x)&=\frac{x^2-1}{x+1} \\
f_3(x)&=\frac{3x-1}{x^3+1} \\
f_4(x)&=\frac{x+1}{x^2-x+1}
\end{aligned} \]
The only such functions are:}
{\( f_2 \), \( f_3 \)}{\( f_1 \), \( f_2 \), \( f_3 \)}{\( f_2 \), \( f_3 \), \( f_4 \)}{\( f_2 \), \( f_4 \)}{}{}}
\LANGpl{
\Question{
Które z podanych funkcji mają punkty nieciągłości 
\[ \begin{aligned}
f_1\colon y&=\frac{x-1}{x^2+1} \\
f_2\colon y&=\frac{x^2-1}{x+1} \\
f_3\colon y&=\frac{3x-1}{x^3+1} \\
f_4\colon y&=\frac{x+1}{x^2-x+1}
\end{aligned} \]
Tylko funkcje:}
{\( f_2 \), \( f_3 \)}{\( f_1 \), \( f_2 \), \( f_3 \)}{\( f_2 \), \( f_3 \), \( f_4 \)}{\( f_2 \), \( f_4 \)}{}{}}
\LANGcs{
\Question{
Rozhodněte, které z následujících funkcí mají body nespojitosti.
\[ \begin{aligned}
f_1\colon y&=\frac{x-1}{x^2+1} \\
f_2\colon y&=\frac{x^2-1}{x+1} \\
f_3\colon y&=\frac{3x-1}{x^3+1} \\
f_4\colon y&=\frac{x+1}{x^2-x+1}
\end{aligned} \]
Jedinými takovými funkcemi jsou:}
{\( f_2 \), \( f_3 \)}{\( f_1 \), \( f_2 \), \( f_3 \)}{\( f_2 \), \( f_3 \), \( f_4 \)}{\( f_2 \), \( f_4 \)}{}{}}
\LANGsk{
\Question{
Rozhodnite, ktoré z nasledujúcich funkcií majú body nespojitosti.
\[ \begin{aligned}
f_1\colon y&=\frac{x-1}{x^2+1} \\
f_2\colon y&=\frac{x^2-1}{x+1} \\
f_3\colon y&=\frac{3x-1}{x^3+1} \\
f_4\colon y&=\frac{x+1}{x^2-x+1}
\end{aligned} \]
Jedinými takými funkciami sú:}
{\( f_2 \), \( f_3 \)}{\( f_1 \), \( f_2 \), \( f_3 \)}{\( f_2 \), \( f_3 \), \( f_4 \)}{\( f_2 \), \( f_4 \)}{}{}}
\LANGes{
\Question{
Decide cuáles de las siguientes funciones tienen puntos de discontinuidad.
\[ \begin{aligned}
f_1(x)&=\frac{x-1}{x^2+1} \\
f_2(x)&=\frac{x^2-1}{x+1} \\
f_3(x)&=\frac{3x-1}{x^3+1} \\
f_4(x)&=\frac{x+1}{x^2-x+1}
\end{aligned} \]
Las únicas funciones así son:}
{\( f_2 \), \( f_3 \)}{\( f_1 \), \( f_2 \), \( f_3 \)}{\( f_2 \), \( f_3 \), \( f_4 \)}{\( f_2 \), \( f_4 \)}{}{}}
\End

\Data\ID{1003085506}
\Updated{2020-07-15 03:07:12}
\Level{B}
\SubArea{Limits and continuity}
\LANGen{
\Question{
Decide which of the following functions have points of discontinuity.
\[ \begin{aligned}
f_1(x)&=\frac1{2^x} \\ 
f_2(x)& =x\cdot3^x \\
f_3(x)&=\frac1{e^x-1} \\
f_4(x)& =e^{2x-1}
\end{aligned} \]
The only such functions are:}
{\( f_3 \)}{\( f_1 \), \( f_3 \)}{\( f_1 \), \( f_4 \)}{\( f_1 \), \( f_3 \), \( f_4 \)}{}{}}
\LANGpl{
\Question{
Które z podanych funkcji mają punkty nieciągłości
\[ \begin{aligned}
f_1\colon y&=\frac1{2^x} \\ 
f_2\colon y& =x\cdot3^x \\
f_3\colon y&=\frac1{e^x-1} \\
f_4\colon y& =e^{2x-1}
\end{aligned} \]
Tylko funkcje:}
{\( f_3 \)}{\( f_1 \), \( f_3 \)}{\( f_1 \), \( f_4 \)}{\( f_1 \), \( f_3 \), \( f_4 \)}{}{}}
\LANGcs{
\Question{
Rozhodněte, které z následujících funkcí mají body nespojitosti.
\[ \begin{aligned}
f_1\colon y&=\frac1{2^x} \\ 
f_2\colon y& =x\cdot3^x \\
f_3\colon y&=\frac1{e^x-1} \\
f_4\colon y& =e^{2x-1}
\end{aligned} \]
Jedinými takovými funkcemi jsou:}
{\( f_3 \)}{\( f_1 \), \( f_3 \)}{\( f_1 \), \( f_4 \)}{\( f_1 \), \( f_3 \), \( f_4 \)}{}{}}
\LANGsk{
\Question{
Rozhodnite, ktoré z nasledujúcich funkcií majú body nespojitosti.
\[ \begin{aligned}
f_1\colon y&=\frac1{2^x} \\ 
f_2\colon y& =x\cdot3^x \\
f_3\colon y&=\frac1{e^x-1} \\
f_4\colon y& =e^{2x-1}
\end{aligned} \]
Jedinými takými funkciami sú:}
{\( f_3 \)}{\( f_1 \), \( f_3 \)}{\( f_1 \), \( f_4 \)}{\( f_1 \), \( f_3 \), \( f_4 \)}{}{}}
\LANGes{
\Question{
Decide cuáles de las siguientes funciones tienen puntos de discontinuidad.
\[ \begin{aligned}
f_1(x)&=\frac1{2^x} \\ 
f_2(x)& =x\cdot3^x \\
f_3(x)&=\frac1{e^x-1} \\
f_4(x)& =e^{2x-1}
\end{aligned} \]
Las únicas funciones así son:}
{\( f_3 \)}{\( f_1 \), \( f_3 \)}{\( f_1 \), \( f_4 \)}{\( f_1 \), \( f_3 \), \( f_4 \)}{}{}}
\End

\Data\ID{1003085507}
\Updated{2021-05-17 08:05:05}
\Level{B}
\SubArea{Limits and continuity}
\LANGen{
\Question{
Decide which of the following functions do not have points of discontinuity.
\[ \begin{aligned}
f_1(x)&=\left\{\begin{array}{lc} x^2 & \text{if } x\leq 1 \\ 2x & \text{ if } x > 1 \end{array}  \right. \\
f_2(x)& =\left\{ \begin{array}{lc} x^2-2x & \text{if } x < -1 \\ 3x & \text{if } x\geq-1 \end{array}  \right. \\
f_3(x)&=\left\{ \begin{array}{lc} 3-x & \text{if } x\leq 2 \\ (x-1)^2 & \text{if } x > 2 \end{array}  \right. \\
f_4(x)&=\left\{ \begin{array}{lc}x^2-2x+1& \text{if } x < 1 \\ \sqrt{x-1} & \text{if } x\geq1 \end{array}  \right.
\end{aligned} \]
The only such functions are:}
{\( f_3 \), \( f_4 \)}{\( f_2 \), \( f_3 \), \( f_4 \)}{\( f_2 \), \( f_3 \)}{\( f_3 \)}{}{}}
\LANGpl{
\Question{
Które z podanych funkcji   nie mają   punktów nieciągłości.
\[ \begin{aligned}
f_1\colon y&=\left\{\begin{array}{lc} x^2 & \text{jeśli } x\leq 1 \\ 2x & \text{jeśli } x > 1 \end{array}  \right. \\
f_2\colon y& =\left\{ \begin{array}{lc} x^2-2x & \text{jeśli } x < -1 \\ 3x & \text{jeśli } x\geq-1 \end{array}  \right. \\
f_3\colon y&=\left\{ \begin{array}{lc} 3-x & \text{jeśli } x\leq 2 \\ (x-1)^2 & \text{jeśli } x > 2 \end{array}  \right. \\
f_4\colon y&=\left\{ \begin{array}{lc}x^2-2x+1& \text{jeśli } x < 1 \\ \sqrt{x-1} & \text{jeśli } x\geq1 \end{array}  \right.
\end{aligned} \]
Tylko funkcje:}
{\( f_3 \), \( f_4 \)}{\( f_2 \), \( f_3 \), \( f_4 \)}{\( f_2 \), \( f_3 \)}{\( f_3 \)}{}{}}
\LANGcs{
\Question{
Rozhodněte, které z následujících funkcí nemají body nespojitosti.
\[ \begin{aligned}
f_1\colon y&=\left\{\begin{array}{lc} x^2 & \text{pro } x\leq 1 \\ 2x & \text{ pro } x > 1 \end{array}  \right. \\
f_2\colon y& =\left\{ \begin{array}{lc} x^2-2x & \text{pro } x < -1 \\ 3x & \text{pro } x\geq-1 \end{array}  \right. \\
f_3\colon y&=\left\{ \begin{array}{lc} 3-x & \text{pro } x\leq 2 \\ (x-1)^2 & \text{pro } x > 2 \end{array}  \right. \\
f_4\colon y&=\left\{ \begin{array}{lc}x^2-2x+1& \text{pro } x < 1 \\ \sqrt{x-1} & \text{pro } x\geq1 \end{array}  \right.
\end{aligned} \]
Jedinými takovými funkcemi jsou:}
{\( f_3 \), \( f_4 \)}{\( f_2 \), \( f_3 \), \( f_4 \)}{\( f_2 \), \( f_3 \)}{\( f_3 \)}{}{}}
\LANGsk{
\Question{
Rozhodnite, ktoré z nasledujúcich funkcií nemajú body nespojitosti.
\[ \begin{aligned}
f_1\colon y&=\left\{\begin{array}{lc} x^2 & \text{pre } x\leq 1 \\ 2x & \text{pre } x > 1 \end{array}  \right. \\
f_2\colon y& =\left\{ \begin{array}{lc} x^2-2x & \text{pre } x < -1 \\ 3x & \text{pre } x\geq-1 \end{array}  \right. \\
f_3\colon y&=\left\{ \begin{array}{lc} 3-x & \text{pre } x\leq 2 \\ (x-1)^2 & \text{pre } x > 2 \end{array}  \right. \\
f_4\colon y&=\left\{ \begin{array}{lc}x^2-2x+1& \text{pre } x < 1 \\ \sqrt{x-1} & \text{pre } x\geq1 \end{array}  \right.
\end{aligned} \]
Jedinými takými funkciami sú:}
{\( f_3 \), \( f_4 \)}{\( f_2 \), \( f_3 \), \( f_4 \)}{\( f_2 \), \( f_3 \)}{\( f_3 \)}{}{}}
\LANGes{
\Question{
Decide cuáles de las siguientes funciones no tienen puntos de discontinuidad.
\[ \begin{aligned}
f_1(x)&=\left\{\begin{array}{lc} x^2 & \text{if } x\leq 1 \\ 2x & \text{ if } x > 1 \end{array}  \right. \\
f_2(x)& =\left\{ \begin{array}{lc} x^2-2x & \text{if } x < -1 \\ 3x & \text{if } x\geq-1 \end{array}  \right. \\
f_3(x)&=\left\{ \begin{array}{lc} 3-x & \text{if } x\leq 2 \\ (x-1)^2 & \text{if } x > 2 \end{array}  \right. \\
f_4(x)&=\left\{ \begin{array}{lc}x^2-2x+1& \text{if } x < 1 \\ \sqrt{x-1} & \text{if } x\geq1 \end{array}  \right.
\end{aligned} \]
Las únicas funciones que lo cumplen son:}
{\( f_3 \), \( f_4 \)}{\( f_2 \), \( f_3 \), \( f_4 \)}{\( f_2 \), \( f_3 \)}{\( f_3 \)}{}{}}
\End

\Data\ID{1003085508}
\Updated{2021-05-17 08:05:29}
\Level{B}
\SubArea{Limits and continuity}
\LANGen{
\Question{
Find all points of discontinuity of the function \( f(x)=\frac{x-3}{x^2-2x-3} \).}
{\( x_1=-1\), \( x_2 = 3 \)}{\( x_1=-1\), \( x_2 = -3 \)}{\( x_1=-1\)}{There is no point of discontinuity.}{}{}}
\LANGpl{
\Question{
Wyznacz wszystkie punkty nieciągłości funkcji \( f\colon y=\frac{x-3}{x^2-2x-3} \).}
{\( x_1=-1\), \( x_2 = 3 \)}{\( x_1=-1\), \( x_2 = -3 \)}{\( x_1=-1\)}{Brak punktów nieciągłości.}{}{}}
\LANGcs{
\Question{
Nalezněte všechny body nespojitosti funkce \( f\colon y=\frac{x-3}{x^2-2x-3} \) .}
{\( x_1=-1\), \( x_2 = 3 \)}{\( x_1=-1\), \( x_2 = -3 \)}{\( x_1=-1\)}{Funkce \(f\) nemá bod nespojitosti.}{}{}}
\LANGsk{
\Question{
Nájdite všetky body nespojitosti funkcie \( f\colon y=\frac{x-3}{x^2-2x-3} \).}
{\( x_1=-1\), \( x_2 = 3 \)}{\( x_1=-1\), \( x_2 = -3 \)}{\( x_1=-1\)}{Funkcia \(f\) nemá bod nespojitosti.}{}{}}
\LANGes{
\Question{
Halla todos los puntos de discontinuidad de la función \( f(x)=\frac{x-3}{x^2-2x-3} \).}
{\( x_1=-1\), \( x_2 = 3 \)}{\( x_1=-1\), \( x_2 = -3 \)}{\( x_1=-1\)}{La función no tiene puntos de discontinuidad.}{}{}}
\End

\Data\ID{1003085509}
\Updated{2021-05-17 08:05:49}
\Level{B}
\SubArea{Limits and continuity}
\LANGen{
\Question{
Find all points of discontinuity of the function \( f(x)=\frac{x^2+4x-5}{x^2+7x+10} \).}
{\( x_1 = -2 \), \( x_2 = -5 \)}{\( x_1 = 2 \), \( x_2 = 5 \)}{\( x_1 = -2 \)}{There is no point of discontinuity.}{}{}}
\LANGpl{
\Question{
Wyznacz wszystkie punkty nieciągłości funkcji \( f\colon y=\frac{x^2+4x-5}{x^2+7x+10} \).}
{\( x_1 = -2 \), \( x_2 = -5 \)}{\( x_1 = 2 \), \( x_2 = 5 \)}{\( x_1 = -2 \)}{Brak punktów nieciągłości.}{}{}}
\LANGcs{
\Question{
Nalezněte všechny body nespojitosti funkce \( f\colon y=\frac{x^2+4x-5}{x^2+7x+10} \).}
{\( x_1 = -2 \), \( x_2 = -5 \)}{\( x_1 = 2 \), \( x_2 = 5 \)}{\( x_1 = -2 \)}{Funkce \(f\) nemá bod nespojitosti.}{}{}}
\LANGsk{
\Question{
Nájdite všetky body nespojitosti funkcie \( f\colon y=\frac{x^2+4x-5}{x^2+7x+10} \).}
{\( x_1 = -2 \), \( x_2 = -5 \)}{\( x_1 = 2 \), \( x_2 = 5 \)}{\( x_1 = -2 \)}{Funkcia \(f\) nemá bod nespojitosti.}{}{}}
\LANGes{
\Question{
Halla todos los puntos de discontinuidad de la función \( f(x)=\frac{x^2+4x-5}{x^2+7x+10} \).}
{\( x_1 = -2 \), \( x_2 = -5 \)}{\( x_1 = 2 \), \( x_2 = 5 \)}{\( x_1 = -2 \)}{La función no tiene puntos de discontinuidad.}{}{}}
\End

\Data\ID{1003085510}
\Updated{2021-05-17 08:05:16}
\Level{B}
\SubArea{Limits and continuity}
\LANGen{
\Question{
Find all points of discontinuity of the function \( f(x)=\frac{16-x^2}{x^2-x+12} \).}
{There is no point of discontinuity.}{\( x_1 = -3 \), \( x_2 = 4 \)}{\( x_1 = -3 \)}{\( x_1 = 3 \), \( x_2 = -4 \)}{}{}}
\LANGpl{
\Question{
Wyznacz wszystkie punkty nieciągłości funkcji \( f\colon y=\frac{16-x^2}{x^2-x+12} \).}
{Brak punktów nieciągłości funkcji.}{\( x_1 = -3 \), \( x_2 = 4 \)}{\( x_1 = -3 \)}{\( x_1 = 3 \), \( x_2 = -4 \)}{}{}}
\LANGcs{
\Question{
Nalezněte všechny body nespojitosti funkce \( f\colon y=\frac{16-x^2}{x^2-x+12} \).}
{Funkce \(f\) nemá bod nespojitosti.}{\( x_1 = -3 \), \( x_2 = 4 \)}{\( x_1 = -3 \)}{\( x_1 = 3 \), \( x_2 = -4 \)}{}{}}
\LANGsk{
\Question{
Nájdite všetky body nespojitosti funkcie \( f\colon y=\frac{16-x^2}{x^2-x+12} \).}
{Funkcia \(f\) nemá bod nespojitosti.}{\( x_1 = -3 \), \( x_2 = 4 \)}{\( x_1 = -3 \)}{\( x_1 = 3 \), \( x_2 = -4 \)}{}{}}
\LANGes{
\Question{
Halla todos los puntos de discontinuidad de la función \( f(x)=\frac{16-x^2}{x^2-x+12} \).}
{La función no tiene  puntos de discontinuidad.}{\( x_1 = -3 \), \( x_2 = 4 \)}{\( x_1 = -3 \)}{\( x_1 = 3 \), \( x_2 = -4 \)}{}{}}
\End

\Data\ID{1003093102}
\Updated{2021-05-17 08:05:18}
\Level{B}
\SubArea{Limits and continuity}
\LANGen{
\Question{
Which of the statements A, B, C, D, E given bellow are incorrect?
\[ \begin{aligned}
\text{A: } & \lim\limits_{x\to-\infty}\left(3-\frac1x\right)=3 \\
\text{B: } & \lim\limits_{x\to-\infty}\left(x^5-2\right)=\infty \\
\text{C: } & \lim\limits_{x\to-\infty}\left(0.3\cdot2^x\right)=-\infty \\
\text{D: } & \lim\limits_{x\to\infty}\left(0.5^x+5\right)= 5 \\
\text{E: } & \lim\limits_{x\to\infty}\left(\log_{\frac12}x-x\right)=0 
\end{aligned} \]
The only incorrect statements are:}
{B, C, E}{B, D}{B, D, E}{A, B, C}{B, C}{}}
\LANGpl{
\Question{
Które z podanych wyrażeń A, B, C, D, E są  niepoprawne?
\[ \begin{aligned}
\text{A: } & \lim\limits_{x\to-\infty}\left(3-\frac1x\right)=3 \\
\text{B: } & \lim\limits_{x\to-\infty}\left(x^5-2\right)=\infty \\
\text{C: } & \lim\limits_{x\to-\infty}\left(0{,}3\cdot2^x\right)=-\infty \\
\text{D: } & \lim\limits_{x\to\infty}\left(0{,}5^x+5\right)= 5 \\
\text{E: } & \lim\limits_{x\to\infty}\left(\log_{\frac12}x-x\right)=0 
\end{aligned} \]
Niepoprawne wyrażenia to:}
{B, C, E}{B, D}{B, D, E}{A, B, C}{B, C}{}}
\LANGcs{
\Question{
Která z níže uvedených tvrzení A, B, C, D, E jsou nepravdivá?
\[ \begin{aligned}
\text{A: } & \lim\limits_{x\to-\infty}\left(3-\frac1x\right)=3 \\
\text{B: } & \lim\limits_{x\to-\infty}\left(x^5-2\right)=\infty \\
\text{C: } & \lim\limits_{x\to-\infty}\left(0{,}3\cdot2^x\right)=-\infty \\
\text{D: } & \lim\limits_{x\to\infty}\left(0{,}5^x+5\right)= 5 \\
\text{E: } & \lim\limits_{x\to\infty}\left(\log_{\frac12}x-x\right)=0 
\end{aligned} \]
Jedinými nepravdivými tvrzeními jsou:}
{B, C, E}{B, D}{B, D, E}{A, B, C}{B, C}{}}
\LANGsk{
\Question{
Ktoré z tvrdení A, B, C, D, E uvedených nižšie nie sú správne?
\[ \begin{aligned}
\text{A: } & \lim\limits_{x\to-\infty}\left(3-\frac1x\right)=3 \\
\text{B: } & \lim\limits_{x\to-\infty}\left(x^5-2\right)=\infty \\
\text{C: } & \lim\limits_{x\to-\infty}\left(0{,}3\cdot2^x\right)=-\infty \\
\text{D: } & \lim\limits_{x\to\infty}\left(0{,}5^x+5\right)= 5 \\
\text{E: } & \lim\limits_{x\to\infty}\left(\log_{\frac12}x-x\right)=0 
\end{aligned} \]
Jedinými nesprávnymi tvrdeniami sú:}
{B, C, E}{B, D}{B, D, E}{A, B, C}{B, C}{}}
\LANGes{
\Question{
¿Cuáles de las siguientes igualdades A, B, C, D, E son falsas?
\[ \begin{aligned}
\text{A: } & \lim\limits_{x\to-\infty}\left(3-\frac1x\right)=3 \\
\text{B: } & \lim\limits_{x\to-\infty}\left(x^5-2\right)=\infty \\
\text{C: } & \lim\limits_{x\to-\infty}\left(0.3\cdot2^x\right)=-\infty \\
\text{D: } & \lim\limits_{x\to\infty}\left(0.5^x+5\right)= 5 \\
\text{E: } & \lim\limits_{x\to\infty}\left(\log_{\frac12}x-x\right)=0 
\end{aligned} \]}
{B, C, E}{B, D}{B, D, E}{A, B, C}{B, C}{}}
\End

\Data\ID{1003109901}
\Updated{2020-07-15 03:07:21}
\Level{B}
\SubArea{Limits and continuity}
\LANGen{
\Question{
Find the next possible correct step in evaluating this limit.
\[ \lim\limits_{x\to\infty}\frac{2x-1}{\sqrt{2x^2-1}} \]}
{\( \lim\limits_{x\to\infty}\frac{2-\frac1x}{\sqrt{2-\frac1{x^2}}} \)}{\( \lim\limits_{x\to\infty}\frac{2-\frac1x}{\sqrt{2x-\frac1x}} \)}{\( \lim\limits_{x\to\infty}\frac{2-\frac1x}{\sqrt{2-x}} \)}{\( \lim\limits_{x\to\infty}\frac{\frac2x-\frac1{x^2}}{\sqrt{2-\frac1{x^2}}} \)}{}{}}
\LANGpl{
\Question{
Wyznacz kolejny możliwy krok do oszacowania tej granicy. 
\[ \lim\limits_{x\to\infty}\frac{2x-1}{\sqrt{2x^2-1}} \]}
{\( \lim\limits_{x\to\infty}\frac{2-\frac1x}{\sqrt{2-\frac1{x^2}}} \)}{\( \lim\limits_{x\to\infty}\frac{2-\frac1x}{\sqrt{2x-\frac1x}} \)}{\( \lim\limits_{x\to\infty}\frac{2-\frac1x}{\sqrt{2-x}} \)}{\( \lim\limits_{x\to\infty}\frac{\frac2x-\frac1{x^2}}{\sqrt{2-\frac1{x^2}}} \)}{}{}}
\LANGcs{
\Question{
Vyberte, jak by mohl po vhodné úpravě pokračovat výpočet limity.
\[ \lim\limits_{x\to\infty}\frac{2x-1}{\sqrt{2x^2-1}} \]}
{\( \lim\limits_{x\to\infty}\frac{2-\frac1x}{\sqrt{2-\frac1{x^2}}} \)}{\( \lim\limits_{x\to\infty}\frac{2-\frac1x}{\sqrt{2x-\frac1x}} \)}{\( \lim\limits_{x\to\infty}\frac{2-\frac1x}{\sqrt{2-x}} \)}{\( \lim\limits_{x\to\infty}\frac{\frac2x-\frac1{x^2}}{\sqrt{2-\frac1{x^2}}} \)}{}{}}
\LANGsk{
\Question{
Vyberte, ako by mohol po vhodnej úprave pokračovať výpočet limity.
\[ \lim\limits_{x\to\infty}\frac{2x-1}{\sqrt{2x^2-1}} \]}
{\( \lim\limits_{x\to\infty}\frac{2-\frac1x}{\sqrt{2-\frac1{x^2}}} \)}{\( \lim\limits_{x\to\infty}\frac{2-\frac1x}{\sqrt{2x-\frac1x}} \)}{\( \lim\limits_{x\to\infty}\frac{2-\frac1x}{\sqrt{2-x}} \)}{\( \lim\limits_{x\to\infty}\frac{\frac2x-\frac1{x^2}}{\sqrt{2-\frac1{x^2}}} \)}{}{}}
\LANGes{
\Question{
Halla el siguiente paso correcto posible para resolver este límite:
\[ \lim\limits_{x\to\infty}\frac{2x-1}{\sqrt{2x^2-1}} \]}
{\( \lim\limits_{x\to\infty}\frac{2-\frac1x}{\sqrt{2-\frac1{x^2}}} \)}{\( \lim\limits_{x\to\infty}\frac{2-\frac1x}{\sqrt{2x-\frac1x}} \)}{\( \lim\limits_{x\to\infty}\frac{2-\frac1x}{\sqrt{2-x}} \)}{\( \lim\limits_{x\to\infty}\frac{\frac2x-\frac1{x^2}}{\sqrt{2-\frac1{x^2}}} \)}{}{}}
\End

\Data\ID{1003109902}
\Updated{2020-07-15 03:07:21}
\Level{B}
\SubArea{Limits and continuity}
\LANGen{
\Question{
Find the next possible correct step in evaluating this limit.
\[ \lim\limits_{x\to\infty}\frac{\sqrt{x^2+1}-x}{x+1} \]}
{\( \lim\limits_{x\to\infty}\frac{\sqrt{1+\frac1{x^2}}-1}{1+\frac1x} \)}{\( \lim\limits_{x\to\infty}\frac{\sqrt{1+\frac1{x^2}}-\frac1x}{1+\frac1x} \)}{\( \lim\limits_{x\to\infty}\frac{\sqrt{x+\frac1x}-1}{1+\frac1x} \)}{\( \lim\limits_{x\to\infty}\frac{\sqrt{1+\frac1{x^2}}-\frac1x}{\frac1x+\frac1{x^2} } \)}{}{}}
\LANGpl{
\Question{
Wyznacz kolejny odpowiedni krok, który należy wykonać, by oszacować granicę.
\[ \lim\limits_{x\to\infty}\frac{\sqrt{x^2+1}-x}{x+1} \]}
{\( \lim\limits_{x\to\infty}\frac{\sqrt{1+\frac1{x^2}}-1}{1+\frac1x} \)}{\( \lim\limits_{x\to\infty}\frac{\sqrt{1+\frac1{x^2}}-\frac1x}{1+\frac1x} \)}{\( \lim\limits_{x\to\infty}\frac{\sqrt{x+\frac1x}-1}{1+\frac1x} \)}{\( \lim\limits_{x\to\infty}\frac{\sqrt{1+\frac1{x^2}}-\frac1x}{\frac1x+\frac1{x^2} } \)}{}{}}
\LANGcs{
\Question{
Vyberte, jak by mohl po vhodné úpravě pokračovat výpočet limity.
\[ \lim\limits_{x\to\infty}\frac{\sqrt{x^2+1}-x}{x+1} \]}
{\( \lim\limits_{x\to\infty}\frac{\sqrt{1+\frac1{x^2}}-1}{1+\frac1x} \)}{\( \lim\limits_{x\to\infty}\frac{\sqrt{1+\frac1{x^2}}-\frac1x}{1+\frac1x} \)}{\( \lim\limits_{x\to\infty}\frac{\sqrt{x+\frac1x}-1}{1+\frac1x} \)}{\( \lim\limits_{x\to\infty}\frac{\sqrt{1+\frac1{x^2}}-\frac1x}{\frac1x+\frac1{x^2} } \)}{}{}}
\LANGsk{
\Question{
Vyberte, ako by mohol po vhodnej úprave pokračovať výpočet limity.
\[ \lim\limits_{x\to\infty}\frac{\sqrt{x^2+1}-x}{x+1} \]}
{\( \lim\limits_{x\to\infty}\frac{\sqrt{1+\frac1{x^2}}-1}{1+\frac1x} \)}{\( \lim\limits_{x\to\infty}\frac{\sqrt{1+\frac1{x^2}}-\frac1x}{1+\frac1x} \)}{\( \lim\limits_{x\to\infty}\frac{\sqrt{x+\frac1x}-1}{1+\frac1x} \)}{\( \lim\limits_{x\to\infty}\frac{\sqrt{1+\frac1{x^2}}-\frac1x}{\frac1x+\frac1{x^2} } \)}{}{}}
\LANGes{
\Question{
Halla el siguiente paso correcto posible para resolver este límite:
\[ \lim\limits_{x\to\infty}\frac{\sqrt{x^2+1}-x}{x+1} \]}
{\( \lim\limits_{x\to\infty}\frac{\sqrt{1+\frac1{x^2}}-1}{1+\frac1x} \)}{\( \lim\limits_{x\to\infty}\frac{\sqrt{1+\frac1{x^2}}-\frac1x}{1+\frac1x} \)}{\( \lim\limits_{x\to\infty}\frac{\sqrt{x+\frac1x}-1}{1+\frac1x} \)}{\( \lim\limits_{x\to\infty}\frac{\sqrt{1+\frac1{x^2}}-\frac1x}{\frac1x+\frac1{x^2} } \)}{}{}}
\End

\Data\ID{1003109903}
\Updated{2020-07-15 03:07:21}
\Level{B}
\SubArea{Limits and continuity}
\LANGen{
\Question{
Find the next possible correct step in evaluating this limit.
\[ \lim\limits_{x\to\infty}\frac{2x^2+3}{\sqrt{3x^4-1}} \]}
{\( \lim\limits_{x\to\infty}\frac{2+\frac3{x^2}}{\sqrt{3-\frac1{x^4}}} \)}{\( \lim\limits_{x\to\infty}\frac{2+\frac3{x^2}}{\sqrt{3x^2-\frac1{x^2}}} \)}{\( \lim\limits_{x\to\infty}\frac{2+\frac3{x^2}}{\sqrt{3x^3-\frac1x}} \)}{\( \lim\limits_{x\to\infty}\frac{\frac2{x^2}+\frac3{x^4}}{\sqrt{3-\frac1{x^4}}} \)}{}{}}
\LANGpl{
\Question{
Wyznacz kolejny możliwy krok, który pomoże w oszacowaniu tej granicy. 
\[ \lim\limits_{x\to\infty}\frac{2x^2+3}{\sqrt{3x^4-1}} \]}
{\( \lim\limits_{x\to\infty}\frac{2+\frac3{x^2}}{\sqrt{3-\frac1{x^4}}} \)}{\( \lim\limits_{x\to\infty}\frac{2+\frac3{x^2}}{\sqrt{3x^2-\frac1{x^2}}} \)}{\( \lim\limits_{x\to\infty}\frac{2+\frac3{x^2}}{\sqrt{3x^3-\frac1x}} \)}{\( \lim\limits_{x\to\infty}\frac{\frac2{x^2}+\frac3{x^4}}{\sqrt{3-\frac1{x^4}}} \)}{}{}}
\LANGcs{
\Question{
Vyberte, jak by mohl po vhodné úpravě pokračovat výpočet limity.
\[ \lim\limits_{x\to\infty}\frac{2x^2+3}{\sqrt{3x^4-1}} \]}
{\( \lim\limits_{x\to\infty}\frac{2+\frac3{x^2}}{\sqrt{3-\frac1{x^4}}} \)}{\( \lim\limits_{x\to\infty}\frac{2+\frac3{x^2}}{\sqrt{3x^2-\frac1{x^2}}} \)}{\( \lim\limits_{x\to\infty}\frac{2+\frac3{x^2}}{\sqrt{3x^3-\frac1x}} \)}{\( \lim\limits_{x\to\infty}\frac{\frac2{x^2}+\frac3{x^4}}{\sqrt{3-\frac1{x^4}}} \)}{}{}}
\LANGsk{
\Question{
Vyberte, ako by mohol po vhodnej úprave pokračovať výpočet limity.
\[ \lim\limits_{x\to\infty}\frac{2x^2+3}{\sqrt{3x^4-1}} \]}
{\( \lim\limits_{x\to\infty}\frac{2+\frac3{x^2}}{\sqrt{3-\frac1{x^4}}} \)}{\( \lim\limits_{x\to\infty}\frac{2+\frac3{x^2}}{\sqrt{3x^2-\frac1{x^2}}} \)}{\( \lim\limits_{x\to\infty}\frac{2+\frac3{x^2}}{\sqrt{3x^3-\frac1x}} \)}{\( \lim\limits_{x\to\infty}\frac{\frac2{x^2}+\frac3{x^4}}{\sqrt{3-\frac1{x^4}}} \)}{}{}}
\LANGes{
\Question{
Halla el siguiente paso correcto posible para resolver este límite:
\[ \lim\limits_{x\to\infty}\frac{2x^2+3}{\sqrt{3x^4-1}} \]}
{\( \lim\limits_{x\to\infty}\frac{2+\frac3{x^2}}{\sqrt{3-\frac1{x^4}}} \)}{\( \lim\limits_{x\to\infty}\frac{2+\frac3{x^2}}{\sqrt{3x^2-\frac1{x^2}}} \)}{\( \lim\limits_{x\to\infty}\frac{2+\frac3{x^2}}{\sqrt{3x^3-\frac1x}} \)}{\( \lim\limits_{x\to\infty}\frac{\frac2{x^2}+\frac3{x^4}}{\sqrt{3-\frac1{x^4}}} \)}{}{}}
\End

\Data\ID{1003109904}
\Updated{2020-07-15 03:07:21}
\Level{B}
\SubArea{Limits and continuity}
\LANGen{
\Question{
Choose the proper expression to expand \( \sqrt{x^2-2}-x \) when evaluating the limit \( \lim\limits_{x\to\infty}\!\left(\sqrt{x^2-2}-x+1\right) \).}
{\( \frac{\sqrt{x^2-2}+x}{\sqrt{x^2-2}+x} \)}{\( \frac{\sqrt{x^2-2}-x}{\sqrt{x^2-2}-x} \)}{\( \frac{\sqrt{x^2-2}}{\sqrt{x^2-2}} \)}{\( \frac{\sqrt{x^2+2}+x}{\sqrt{x^2+2}+x} \)}{}{}}
\LANGpl{
\Question{
Wybierz odpowiednie wyrażenie, aby rozszerzyć \( \sqrt{x^2-2}-x \) wyznaczając granicę  \( \lim\limits_{x\to\infty}\!\left(\sqrt{x^2-2}-x+1\right) \).}
{\( \frac{\sqrt{x^2-2}+x}{\sqrt{x^2-2}+x} \)}{\( \frac{\sqrt{x^2-2}-x}{\sqrt{x^2-2}-x} \)}{\( \frac{\sqrt{x^2-2}}{\sqrt{x^2-2}} \)}{\( \frac{\sqrt{x^2+2}+x}{\sqrt{x^2+2}+x} \)}{}{}}
\LANGcs{
\Question{
Vyberte, kterým z následujících výrazů je vhodné rozšířit výraz \( \sqrt{x^2-2}-x \) při výpočtu limity \( \lim\limits_{x\to\infty}\!\left(\sqrt{x^2-2}-x+1\right) \).}
{\( \frac{\sqrt{x^2-2}+x}{\sqrt{x^2-2}+x} \)}{\( \frac{\sqrt{x^2-2}-x}{\sqrt{x^2-2}-x} \)}{\( \frac{\sqrt{x^2-2}}{\sqrt{x^2-2}} \)}{\( \frac{\sqrt{x^2+2}+x}{\sqrt{x^2+2}+x} \)}{}{}}
\LANGsk{
\Question{
Vyberte, ktorým z následujúcich výrazov je vhodné rozšíriť výraz \( \sqrt{x^2-2}-x \) pri výpočte limity \( \lim\limits_{x\to\infty}\!\left(\sqrt{x^2-2}-x+1\right) \).}
{\( \frac{\sqrt{x^2-2}+x}{\sqrt{x^2-2}+x} \)}{\( \frac{\sqrt{x^2-2}-x}{\sqrt{x^2-2}-x} \)}{\( \frac{\sqrt{x^2-2}}{\sqrt{x^2-2}} \)}{\( \frac{\sqrt{x^2+2}+x}{\sqrt{x^2+2}+x} \)}{}{}}
\LANGes{
\Question{
Elige la expresión adecuada para expandir \( \sqrt{x^2-2}-x \) al resolver el límite \( \lim\limits_{x\to\infty}\!\left(\sqrt{x^2-2}-x+1\right) \).}
{\( \frac{\sqrt{x^2-2}+x}{\sqrt{x^2-2}+x} \)}{\( \frac{\sqrt{x^2-2}-x}{\sqrt{x^2-2}-x} \)}{\( \frac{\sqrt{x^2-2}}{\sqrt{x^2-2}} \)}{\( \frac{\sqrt{x^2+2}+x}{\sqrt{x^2+2}+x} \)}{}{}}
\End

\Data\ID{1003109905}
\Updated{2020-07-15 03:07:21}
\Level{B}
\SubArea{Limits and continuity}
\LANGen{
\Question{
Choose the proper expression to expand \( \sqrt{x-5}-\sqrt x \)  when evaluating the limit \( \lim\limits_{x\to\infty}\!\left(\sqrt{x-5}-\sqrt x-1 \right) \).}
{\( \frac{\sqrt{x-5}+\sqrt x}{\sqrt{x-5}+\sqrt x} \)}{\( \frac{\sqrt{x-5}+\sqrt x+1}{\sqrt{x-5}+\sqrt x+1 } \)}{\( \frac{\sqrt{x+5}+\sqrt x}{\sqrt{x+5}+\sqrt x} \)}{\( \frac{\sqrt{x-5}}{\sqrt{x-5}} \)}{\( \frac{\sqrt{x-5}+\sqrt x-1}{\sqrt{x-5}+\sqrt x-1} \)}{}}
\LANGpl{
\Question{
Wybierz odpowiednie wyrażenie, aby rozszerzyć  \( \sqrt{x-5}-\sqrt x \)  wyznaczając granicę \( \lim\limits_{x\to\infty}\!\left(\sqrt{x-5}-\sqrt x-1 \right) \).}
{\( \frac{\sqrt{x-5}+\sqrt x}{\sqrt{x-5}+\sqrt x} \)}{\( \frac{\sqrt{x-5}+\sqrt x+1}{\sqrt{x-5}+\sqrt x+1 } \)}{\( \frac{\sqrt{x+5}+\sqrt x}{\sqrt{x+5}+\sqrt x} \)}{\( \frac{\sqrt{x-5}}{\sqrt{x-5}} \)}{\( \frac{\sqrt{x-5}+\sqrt x-1}{\sqrt{x-5}+\sqrt x-1} \)}{}}
\LANGcs{
\Question{
Vyberte, kterým z následujících výrazů je vhodné rozšířit výraz \( \sqrt{x-5}-\sqrt x \)  při výpočtu limity \( \lim\limits_{x\to\infty}\!\left(\sqrt{x-5}-\sqrt x-1 \right) \).}
{\( \frac{\sqrt{x-5}+\sqrt x}{\sqrt{x-5}+\sqrt x} \)}{\( \frac{\sqrt{x-5}+\sqrt x+1}{\sqrt{x-5}+\sqrt x+1 } \)}{\( \frac{\sqrt{x+5}+\sqrt x}{\sqrt{x+5}+\sqrt x} \)}{\( \frac{\sqrt{x-5}}{\sqrt{x-5}} \)}{\( \frac{\sqrt{x-5}+\sqrt x-1}{\sqrt{x-5}+\sqrt x-1} \)}{}}
\LANGsk{
\Question{
Vyberte, ktorým z následujúcich výrazov je vhodné rozšíriť výraz \( \sqrt{x-5}-\sqrt x \)  pri výpočte limity \( \lim\limits_{x\to\infty}\!\left(\sqrt{x-5}-\sqrt x-1 \right) \).}
{\( \frac{\sqrt{x-5}+\sqrt x}{\sqrt{x-5}+\sqrt x} \)}{\( \frac{\sqrt{x-5}+\sqrt x+1}{\sqrt{x-5}+\sqrt x+1 } \)}{\( \frac{\sqrt{x+5}+\sqrt x}{\sqrt{x+5}+\sqrt x} \)}{\( \frac{\sqrt{x-5}}{\sqrt{x-5}} \)}{\( \frac{\sqrt{x-5}+\sqrt x-1}{\sqrt{x-5}+\sqrt x-1} \)}{}}
\LANGes{
\Question{
Elige la expresión adecuada para expandir \( \sqrt{x-5}-\sqrt x \)  al resolver el límite \( \lim\limits_{x\to\infty}\!\left(\sqrt{x-5}-\sqrt x-1 \right) \).}
{\( \frac{\sqrt{x-5}+\sqrt x}{\sqrt{x-5}+\sqrt x} \)}{\( \frac{\sqrt{x-5}+\sqrt x+1}{\sqrt{x-5}+\sqrt x+1 } \)}{\( \frac{\sqrt{x+5}+\sqrt x}{\sqrt{x+5}+\sqrt x} \)}{\( \frac{\sqrt{x-5}}{\sqrt{x-5}} \)}{\( \frac{\sqrt{x-5}+\sqrt x-1}{\sqrt{x-5}+\sqrt x-1} \)}{}}
\End

\Data\ID{1003109906}
\Updated{2020-07-15 03:07:21}
\Level{B}
\SubArea{Limits and continuity}
\LANGen{
\Question{
Evaluate the limit.
\[ \lim\limits_{x\to\infty}\frac{\sqrt{3x^2-1}}{x+2} \]}
{\( \sqrt3 \)}{\( \frac{\sqrt3}2 \)}{\( 1 \)}{\( \frac{\sqrt3}3 \)}{}{}}
\LANGpl{
\Question{
Oszacuj granicę. 
\[ \lim\limits_{x\to\infty}\frac{\sqrt{3x^2-1}}{x+2} \]}
{\( \sqrt3 \)}{\( \frac{\sqrt3}2 \)}{\( 1 \)}{\( \frac{\sqrt3}3 \)}{}{}}
\LANGcs{
\Question{
Vypočítejte následující limitu.
\[ \lim\limits_{x\to\infty}\frac{\sqrt{3x^2-1}}{x+2} \]}
{\( \sqrt3 \)}{\( \frac{\sqrt3}2 \)}{\( 1 \)}{\( \frac{\sqrt3}3 \)}{}{}}
\LANGsk{
\Question{
Vypočítajte následujúcu limitu.
\[ \lim\limits_{x\to\infty}\frac{\sqrt{3x^2-1}}{x+2} \]}
{\( \sqrt3 \)}{\( \frac{\sqrt3}2 \)}{\( 1 \)}{\( \frac{\sqrt3}3 \)}{}{}}
\LANGes{
\Question{
Resuelve el siguiente límite:
\[ \lim\limits_{x\to\infty}\frac{\sqrt{3x^2-1}}{x+2} \]}
{\( \sqrt3 \)}{\( \frac{\sqrt3}2 \)}{\( 1 \)}{\( \frac{\sqrt3}3 \)}{}{}}
\End

\Data\ID{1003109907}
\Updated{2020-07-15 03:07:21}
\Level{B}
\SubArea{Limits and continuity}
\LANGen{
\Question{
Evaluate the limit.
\[ \lim\limits_{x\to\infty}\frac{4x-\sqrt x}{2-x} \]}
{\( -4 \)}{\( 4 \)}{\( 2 \)}{\( -3 \)}{}{}}
\LANGpl{
\Question{
Wyznacz granicę.
\[ \lim\limits_{x\to\infty}\frac{4x-\sqrt x}{2-x} \]}
{\( -4 \)}{\( 4 \)}{\( 2 \)}{\( -3 \)}{}{}}
\LANGcs{
\Question{
Vypočítejte následující limitu.
\[ \lim\limits_{x\to\infty}\frac{4x-\sqrt x}{2-x} \]}
{\( -4 \)}{\( 4 \)}{\( 2 \)}{\( -3 \)}{}{}}
\LANGsk{
\Question{
Vypočítajte nasledujúcu limitu.
\[ \lim\limits_{x\to\infty}\frac{4x-\sqrt x}{2-x} \]}
{\( -4 \)}{\( 4 \)}{\( 2 \)}{\( -3 \)}{}{}}
\LANGes{
\Question{
Resuelve el siguiente límite:
\[ \lim\limits_{x\to\infty}\frac{4x-\sqrt x}{2-x} \]}
{\( -4 \)}{\( 4 \)}{\( 2 \)}{\( -3 \)}{}{}}
\End

\Data\ID{1003109908}
\Updated{2020-07-15 03:07:21}
\Level{B}
\SubArea{Limits and continuity}
\LANGen{
\Question{
Evaluate the limit.
\[ \lim\limits_{x\to\infty}\frac{3x+1}{\sqrt{4x^2-1}-x} \]}
{\( 3 \)}{\( \frac32 \)}{\( 1 \)}{\( 0 \)}{\( \infty \)}{}}
\LANGpl{
\Question{
Wyznacz granicę.
\[ \lim\limits_{x\to\infty}\frac{3x+1}{\sqrt{4x^2-1}-x} \]}
{\( 3 \)}{\( \frac32 \)}{\( 1 \)}{\( 0 \)}{\( \infty \)}{}}
\LANGcs{
\Question{
Vypočítejte následující limitu.
\[ \lim\limits_{x\to\infty}\frac{3x+1}{\sqrt{4x^2-1}-x} \]}
{\( 3 \)}{\( \frac32 \)}{\( 1 \)}{\( 0 \)}{\( \infty \)}{}}
\LANGsk{
\Question{
Vypočítajte následujúcu limitu.
\[ \lim\limits_{x\to\infty}\frac{3x+1}{\sqrt{4x^2-1}-x} \]}
{\( 3 \)}{\( \frac32 \)}{\( 1 \)}{\( 0 \)}{\( \infty \)}{}}
\LANGes{
\Question{
Resuelve el siguiente límite:
\[ \lim\limits_{x\to\infty}\frac{3x+1}{\sqrt{4x^2-1}-x} \]}
{\( 3 \)}{\( \frac32 \)}{\( 1 \)}{\( 0 \)}{\( \infty \)}{}}
\End

\Data\ID{1003109909}
\Updated{2020-07-15 03:07:21}
\Level{B}
\SubArea{Limits and continuity}
\LANGen{
\Question{
Evaluate the limit.
\[ \lim\limits_{x\to\infty}\!\left(\sqrt{x-3}-\sqrt x\right) \]}
{\( 0 \)}{\( \infty \)}{\( -\infty \)}{\( -3 \)}{}{}}
\LANGpl{
\Question{
Wyznacz granicę. 
\[ \lim\limits_{x\to\infty}\!\left(\sqrt{x-3}-\sqrt x\right) \]}
{\( 0 \)}{\( \infty \)}{\( -\infty \)}{\( -3 \)}{}{}}
\LANGcs{
\Question{
Vypočítejte následující limitu.
\[ \lim\limits_{x\to\infty}\!\left(\sqrt{x-3}-\sqrt x\right) \]}
{\( 0 \)}{\( \infty \)}{\( -\infty \)}{\( -3 \)}{}{}}
\LANGsk{
\Question{
Vypočítajte následujúcu limitu.
\[ \lim\limits_{x\to\infty}\!\left(\sqrt{x-3}-\sqrt x\right) \]}
{\( 0 \)}{\( \infty \)}{\( -\infty \)}{\( -3 \)}{}{}}
\LANGes{
\Question{
Resuelve el siguiente límite:
\[ \lim\limits_{x\to\infty}\!\left(\sqrt{x-3}-\sqrt x\right) \]}
{\( 0 \)}{\( \infty \)}{\( -\infty \)}{\( -3 \)}{}{}}
\End

\Data\ID{1003109910}
\Updated{2020-07-15 03:07:21}
\Level{B}
\SubArea{Limits and continuity}
\LANGen{
\Question{
Evaluate the limit.
\[ \lim\limits_{x\to-\infty}\!\left(2+x+\sqrt{x^2-1}\right) \]}
{\( 2 \)}{\( \infty \)}{\( -\infty \)}{\( 0 \)}{}{}}
\LANGpl{
\Question{
Wyznacz granicę.
\[ \lim\limits_{x\to-\infty}\!\left(2+x+\sqrt{x^2-1}\right) \]}
{\( 2 \)}{\( \infty \)}{\( -\infty \)}{\( 0 \)}{}{}}
\LANGcs{
\Question{
Vypočítejte následující limitu.
\[ \lim\limits_{x\to-\infty}\!\left(2+x+\sqrt{x^2-1}\right) \]}
{\( 2 \)}{\( \infty \)}{\( -\infty \)}{\( 0 \)}{}{}}
\LANGsk{
\Question{
Vypočítajte nasledujúcu limitu.
\[ \lim\limits_{x\to-\infty}\!\left(2+x+\sqrt{x^2-1}\right) \]}
{\( 2 \)}{\( \infty \)}{\( -\infty \)}{\( 0 \)}{}{}}
\LANGes{
\Question{
Resuelve el siguiente límite:
\[ \lim\limits_{x\to-\infty}\!\left(2+x+\sqrt{x^2-1}\right) \]}
{\( 2 \)}{\( \infty \)}{\( -\infty \)}{\( 0 \)}{}{}}
\End

\Data\ID{1003109911}
\Updated{2020-07-15 04:07:40}
\Level{B}
\SubArea{Limits and continuity}
\LANGen{
\Question{
Evaluate the limit.
\[ \lim\limits_{x\to\infty}\!\left(\sqrt{2x+9x^2}-3x\right) \]}
{\( \frac13 \)}{\( \frac23 \)}{\( 0 \)}{\( \infty \)}{}{}}
\LANGpl{
\Question{
Wyznacz granicę.
\[ \lim\limits_{x\to\infty}\!\left(\sqrt{2x+9x^2}-3x\right) \]}
{\( \frac13 \)}{\( \frac23 \)}{\( 0 \)}{\( \infty \)}{}{}}
\LANGcs{
\Question{
Vypočítejte následující limitu.
\[ \lim\limits_{x\to\infty}\!\left(\sqrt{2x+9x^2}-3x\right) \]}
{\( \frac13 \)}{\( \frac23 \)}{\( 0 \)}{\( \infty \)}{}{}}
\LANGsk{
\Question{
Vypočítajte nasledujúcu limitu.
\[ \lim\limits_{x\to\infty}\!\left(\sqrt{2x+9x^2}-3x\right) \]}
{\( \frac13 \)}{\( \frac23 \)}{\( 0 \)}{\( \infty \)}{}{}}
\LANGes{
\Question{
Resuelve el siguiente límite:
\[ \lim\limits_{x\to\infty}\!\left(\sqrt{2x+9x^2}-3x\right) \]}
{\( \frac13 \)}{\( \frac23 \)}{\( 0 \)}{\( \infty \)}{}{}}
\End

\Data\ID{1003109912}
\Updated{2020-07-15 03:07:21}
\Level{B}
\SubArea{Limits and continuity}
\LANGen{
\Question{
Evaluate the limit.
\[ \lim\limits_{x\to\infty}\!\left( x-\sqrt{x^2-x+1} \right) \]}
{\( \frac12 \)}{\( 1 \)}{\( 0 \)}{\( -\infty \)}{\( -\frac12 \)}{}}
\LANGpl{
\Question{
Wyznacz granicę.
\[ \lim\limits_{x\to\infty}\!\left( x-\sqrt{x^2-x+1} \right) \]}
{\( \frac12 \)}{\( 1 \)}{\( 0 \)}{\( -\infty \)}{\( -\frac12 \)}{}}
\LANGcs{
\Question{
Vypočítejte následující limitu.
\[ \lim\limits_{x\to\infty}\!\left( x-\sqrt{x^2-x+1} \right) \]}
{\( \frac12 \)}{\( 1 \)}{\( 0 \)}{\( -\infty \)}{\( -\frac12 \)}{}}
\LANGsk{
\Question{
Vypočítajte nasledujúcu limitu.
\[ \lim\limits_{x\to\infty}\!\left( x-\sqrt{x^2-x+1} \right) \]}
{\( \frac12 \)}{\( 1 \)}{\( 0 \)}{\( -\infty \)}{\( -\frac12 \)}{}}
\LANGes{
\Question{
Resuelve el siguiente límite:
\[ \lim\limits_{x\to\infty}\!\left( x-\sqrt{x^2-x+1} \right) \]}
{\( \frac12 \)}{\( 1 \)}{\( 0 \)}{\( -\infty \)}{\( -\frac12 \)}{}}
\End

\Data\ID{1003164301}
\Updated{2020-07-15 03:07:21}
\Level{C}
\SubArea{Limits and continuity}
\LANGen{
\Question{
Which of the following situations could arise for suitable functions \( f \) and \( g \)?}
{\( \lim\limits_{x\to3} f(x)=\infty\ \wedge\ \lim\limits_{x\to3} g(x)=\infty\ \wedge\ \lim\limits_{x\to3}\frac{f(x)}{g(x)}=5 \)}{\( \lim\limits_{x\to3} f(x)=1\ \wedge\ \lim\limits_{x\to3} g(x)=\infty\ \wedge\ \lim\limits_{x\to3}\frac{f(x)}{g(x)}=5 \)}{\( \lim\limits_{x\to3} f(x)=\infty\ \wedge\ \lim\limits_{x\to3} g(x)=1\ \wedge\ \lim\limits_{x\to3}\frac{f(x)}{g(x)}=5 \)}{\( \lim\limits_{x\to3} f(x)=0\ \wedge\ \lim\limits_{x\to3} g(x)=\infty\ \wedge\ \lim\limits_{x\to3}\frac{f(x)}{g(x)}=5 \)}{}{}}
\LANGpl{
\Question{
Które z poniższych sytuacji mogą zachodzić dla odpowiednich funkcji \( f \) i \( g \)?}
{\( \lim\limits_{x\to3} f(x)=\infty\ \wedge\ \lim\limits_{x\to3} g(x)=\infty\ \wedge\ \lim\limits_{x\to3}\frac{f(x)}{g(x)}=5 \)}{\( \lim\limits_{x\to3} f(x)=1\ \wedge\ \lim\limits_{x\to3} g(x)=\infty\ \wedge\ \lim\limits_{x\to3}\frac{f(x)}{g(x)}=5 \)}{\( \lim\limits_{x\to3} f(x)=\infty\ \wedge\ \lim\limits_{x\to3} g(x)=1\ \wedge\ \lim\limits_{x\to3}\frac{f(x)}{g(x)}=5 \)}{\( \lim\limits_{x\to3} f(x)=0\ \wedge\ \lim\limits_{x\to3} g(x)=\infty\ \wedge\ \lim\limits_{x\to3}\frac{f(x)}{g(x)}=5 \)}{}{}}
\LANGcs{
\Question{
Která z následujících situací může pro vhodné funkce \( f \) a \( g \) nastat?}
{\( \lim\limits_{x\to3} f(x)=\infty\ \wedge\ \lim\limits_{x\to3} g(x)=\infty\ \wedge\ \lim\limits_{x\to3}\frac{f(x)}{g(x)}=5 \)}{\( \lim\limits_{x\to3} f(x)=1\ \wedge\ \lim\limits_{x\to3} g(x)=\infty\ \wedge\ \lim\limits_{x\to3}\frac{f(x)}{g(x)}=5 \)}{\( \lim\limits_{x\to3} f(x)=\infty\ \wedge\ \lim\limits_{x\to3} g(x)=1\ \wedge\ \lim\limits_{x\to3}\frac{f(x)}{g(x)}=5 \)}{\( \lim\limits_{x\to3} f(x)=0\ \wedge\ \lim\limits_{x\to3} g(x)=\infty\ \wedge\ \lim\limits_{x\to3}\frac{f(x)}{g(x)}=5 \)}{}{}}
\LANGsk{
\Question{
Ktorá z nasledujúcich situácií môže pri vhodnej funkcii \( f \) a \( g \) nastať?}
{\( \lim\limits_{x\to3} f(x)=\infty\ \wedge\ \lim\limits_{x\to3} g(x)=\infty\ \wedge\ \lim\limits_{x\to3}\frac{f(x)}{g(x)}=5 \)}{\( \lim\limits_{x\to3} f(x)=1\ \wedge\ \lim\limits_{x\to3} g(x)=\infty\ \wedge\ \lim\limits_{x\to3}\frac{f(x)}{g(x)}=5 \)}{\( \lim\limits_{x\to3} f(x)=\infty\ \wedge\ \lim\limits_{x\to3} g(x)=1\ \wedge\ \lim\limits_{x\to3}\frac{f(x)}{g(x)}=5 \)}{\( \lim\limits_{x\to3} f(x)=0\ \wedge\ \lim\limits_{x\to3} g(x)=\infty\ \wedge\ \lim\limits_{x\to3}\frac{f(x)}{g(x)}=5 \)}{}{}}
\LANGes{
\Question{
¿Cuál de las siguientes situaciones podría surgir para las funciones adecuadas \( f \) y \( g \)?}
{\( \lim\limits_{x\to3} f(x)=\infty\ \wedge\ \lim\limits_{x\to3} g(x)=\infty\ \wedge\ \lim\limits_{x\to3}\frac{f(x)}{g(x)}=5 \)}{\( \lim\limits_{x\to3} f(x)=1\ \wedge\ \lim\limits_{x\to3} g(x)=\infty\ \wedge\ \lim\limits_{x\to3}\frac{f(x)}{g(x)}=5 \)}{\( \lim\limits_{x\to3} f(x)=\infty\ \wedge\ \lim\limits_{x\to3} g(x)=1\ \wedge\ \lim\limits_{x\to3}\frac{f(x)}{g(x)}=5 \)}{\( \lim\limits_{x\to3} f(x)=0\ \wedge\ \lim\limits_{x\to3} g(x)=\infty\ \wedge\ \lim\limits_{x\to3}\frac{f(x)}{g(x)}=5 \)}{}{}}
\End

\Data\ID{1003164302}
\Updated{2020-07-15 03:07:21}
\Level{C}
\SubArea{Limits and continuity}
\LANGen{
\Question{
Which of the following situations could arise for suitable functions \( f \) and \( g \)?}
{\( \lim\limits_{x\to2}f(x)=\infty\ \wedge\ \lim\limits_{x\to2}g(x)=\infty\ \wedge\ \lim\limits_{x\to2}[f(x)-g(x)]=\infty \)}{\( \lim\limits_{x\to2}f(x)=1\ \wedge\ \lim\limits_{x\to2}g(x)=\infty\ \wedge\ \lim\limits_{x\to2}\frac{f(x)}{g(x)}=1 \)}{\( \lim\limits_{x\to2}f(x)=-\infty\ \wedge\ \lim\limits_{x\to2}g(x)=1\ \wedge\ \lim\limits_{x\to2}[f(x)+g(x)]=1 \)}{\( \lim\limits_{x\to2}f(x)=-\infty\ \wedge\ \lim\limits_{x\to2}g(x)=-\infty\ \wedge\ \lim\limits_{x\to2}[f(x)\cdot g(x)]=-\infty \)}{}{}}
\LANGpl{
\Question{
Które z poniższych sytuacji mogą zachodzić dla odpowiednich funkcji \( f \) i \( g \)?}
{\( \lim\limits_{x\to2}f(x)=\infty\ \wedge\ \lim\limits_{x\to2}g(x)=\infty\ \wedge\ \lim\limits_{x\to2}[f(x)-g(x)]=\infty \)}{\( \lim\limits_{x\to2}f(x)=1\ \wedge\ \lim\limits_{x\to2}g(x)=\infty\ \wedge\ \lim\limits_{x\to2}\frac{f(x)}{g(x)}=1 \)}{\( \lim\limits_{x\to2}f(x)=-\infty\ \wedge\ \lim\limits_{x\to2}g(x)=1\ \wedge\ \lim\limits_{x\to2}[f(x)+g(x)]=1 \)}{\( \lim\limits_{x\to2}f(x)=-\infty\ \wedge\ \lim\limits_{x\to2}g(x)=-\infty\ \wedge\ \lim\limits_{x\to2}[f(x)\cdot g(x)]=-\infty \)}{}{}}
\LANGcs{
\Question{
Která z následujících situací může pro vhodné funkce \( f \) a \( g \) nastat?}
{\( \lim\limits_{x\to2}f(x)=\infty\ \wedge\ \lim\limits_{x\to2}g(x)=\infty\ \wedge\ \lim\limits_{x\to2}[f(x)-g(x)]=\infty \)}{\( \lim\limits_{x\to2}f(x)=1\ \wedge\ \lim\limits_{x\to2}g(x)=\infty\ \wedge\ \lim\limits_{x\to2}\frac{f(x)}{g(x)}=1 \)}{\( \lim\limits_{x\to2}f(x)=-\infty\ \wedge\ \lim\limits_{x\to2}g(x)=1\ \wedge\ \lim\limits_{x\to2}[f(x)+g(x)]=1 \)}{\( \lim\limits_{x\to2}f(x)=-\infty\ \wedge\ \lim\limits_{x\to2}g(x)=-\infty\ \wedge\ \lim\limits_{x\to2}[f(x)\cdot g(x)]=-\infty \)}{}{}}
\LANGsk{
\Question{
Ktorá z následujúcich situácií môže pri vhodnej funkcii \( f \) a \( g \) nastať?}
{\( \lim\limits_{x\to2}f(x)=\infty\ \wedge\ \lim\limits_{x\to2}g(x)=\infty\ \wedge\ \lim\limits_{x\to2}[f(x)-g(x)]=\infty \)}{\( \lim\limits_{x\to2}f(x)=1\ \wedge\ \lim\limits_{x\to2}g(x)=\infty\ \wedge\ \lim\limits_{x\to2}\frac{f(x)}{g(x)}=1 \)}{\( \lim\limits_{x\to2}f(x)=-\infty\ \wedge\ \lim\limits_{x\to2}g(x)=1\ \wedge\ \lim\limits_{x\to2}[f(x)+g(x)]=1 \)}{\( \lim\limits_{x\to2}f(x)=-\infty\ \wedge\ \lim\limits_{x\to2}g(x)=-\infty\ \wedge\ \lim\limits_{x\to2}[f(x)\cdot g(x)]=-\infty \)}{}{}}
\LANGes{
\Question{
¿Cuál de las siguientes situaciones podría surgir para las funciones adecuadas \( f \) y \( g \)?}
{\( \lim\limits_{x\to2}f(x)=\infty\ \wedge\ \lim\limits_{x\to2}g(x)=\infty\ \wedge\ \lim\limits_{x\to2}[f(x)-g(x)]=\infty \)}{\( \lim\limits_{x\to2}f(x)=1\ \wedge\ \lim\limits_{x\to2}g(x)=\infty\ \wedge\ \lim\limits_{x\to2}\frac{f(x)}{g(x)}=1 \)}{\( \lim\limits_{x\to2}f(x)=-\infty\ \wedge\ \lim\limits_{x\to2}g(x)=1\ \wedge\ \lim\limits_{x\to2}[f(x)+g(x)]=1 \)}{\( \lim\limits_{x\to2}f(x)=-\infty\ \wedge\ \lim\limits_{x\to2}g(x)=-\infty\ \wedge\ \lim\limits_{x\to2}[f(x)\cdot g(x)]=-\infty \)}{}{}}
\End

\Data\ID{1003164303}
\Updated{2020-07-15 03:07:21}
\Level{C}
\SubArea{Limits and continuity}
\LANGen{
\Question{
Which of the following situations could arise for suitable functions \( f \) and \( g \)?}
{\( \lim\limits_{x\to5}f(x)=0\ \wedge\ \lim\limits_{x\to5}g(x)=\infty\ \wedge\ \lim\limits_{x\to5}[f(x)\cdot g(x)]=13 \)}{\( \lim\limits_{x\to5}f(x)=1\ \wedge\ \lim\limits_{x\to5}g(x)=\infty\ \wedge\ \lim\limits_{x\to5}[f(x)\cdot g(x)]=13 \)}{\( \lim\limits_{x\to5}f(x)=\infty\ \wedge\ \lim\limits_{x\to5}g(x)=\infty\ \wedge\ \lim\limits_{x\to5}[f(x)\cdot g(x)]=13 \)}{\( \lim\limits_{x\to5}f(x)=-\infty\ \wedge\ \lim\limits_{x\to5}g(x)=\infty\ \wedge\ \lim\limits_{x\to5}[f(x)\cdot g(x)]=13 \)}{}{}}
\LANGpl{
\Question{
Które z poniższych sytuacji mogą zachodzić dla odpowiednich funkcji \( f \) i \( g \)?}
{\( \lim\limits_{x\to5}f(x)=0\ \wedge\ \lim\limits_{x\to5}g(x)=\infty\ \wedge\ \lim\limits_{x\to5}[f(x)\cdot g(x)]=13 \)}{\( \lim\limits_{x\to5}f(x)=1\ \wedge\ \lim\limits_{x\to5}g(x)=\infty\ \wedge\ \lim\limits_{x\to5}[f(x)\cdot g(x)]=13 \)}{\( \lim\limits_{x\to5}f(x)=\infty\ \wedge\ \lim\limits_{x\to5}g(x)=\infty\ \wedge\ \lim\limits_{x\to5}[f(x)\cdot g(x)]=13 \)}{\( \lim\limits_{x\to5}f(x)=-\infty\ \wedge\ \lim\limits_{x\to5}g(x)=\infty\ \wedge\ \lim\limits_{x\to5}[f(x)\cdot g(x)]=13 \)}{}{}}
\LANGcs{
\Question{
Která z následujících situací může pro vhodné funkce \( f \) a \( g \) nastat?}
{\( \lim\limits_{x\to5}f(x)=0\ \wedge\ \lim\limits_{x\to5}g(x)=\infty\ \wedge\ \lim\limits_{x\to5}[f(x)\cdot g(x)]=13 \)}{\( \lim\limits_{x\to5}f(x)=1\ \wedge\ \lim\limits_{x\to5}g(x)=\infty\ \wedge\ \lim\limits_{x\to5}[f(x)\cdot g(x)]=13 \)}{\( \lim\limits_{x\to5}f(x)=\infty\ \wedge\ \lim\limits_{x\to5}g(x)=\infty\ \wedge\ \lim\limits_{x\to5}[f(x)\cdot g(x)]=13 \)}{\( \lim\limits_{x\to5}f(x)=-\infty\ \wedge\ \lim\limits_{x\to5}g(x)=\infty\ \wedge\ \lim\limits_{x\to5}[f(x)\cdot g(x)]=13 \)}{}{}}
\LANGsk{
\Question{
Ktorá z nasledujúcich situácií môže pri vhodnej funkcii \( f \) a \( g \) nastať?}
{\( \lim\limits_{x\to5}f(x)=0\ \wedge\ \lim\limits_{x\to5}g(x)=\infty\ \wedge\ \lim\limits_{x\to5}[f(x)\cdot g(x)]=13 \)}{\( \lim\limits_{x\to5}f(x)=1\ \wedge\ \lim\limits_{x\to5}g(x)=\infty\ \wedge\ \lim\limits_{x\to5}[f(x)\cdot g(x)]=13 \)}{\( \lim\limits_{x\to5}f(x)=\infty\ \wedge\ \lim\limits_{x\to5}g(x)=\infty\ \wedge\ \lim\limits_{x\to5}[f(x)\cdot g(x)]=13 \)}{\( \lim\limits_{x\to5}f(x)=-\infty\ \wedge\ \lim\limits_{x\to5}g(x)=\infty\ \wedge\ \lim\limits_{x\to5}[f(x)\cdot g(x)]=13 \)}{}{}}
\LANGes{
\Question{
¿Cuál de las siguientes situaciones podría surgir para las funciones adecuadas \( f \) y \( g \)?}
{\( \lim\limits_{x\to5}f(x)=0\ \wedge\ \lim\limits_{x\to5}g(x)=\infty\ \wedge\ \lim\limits_{x\to5}[f(x)\cdot g(x)]=13 \)}{\( \lim\limits_{x\to5}f(x)=1\ \wedge\ \lim\limits_{x\to5}g(x)=\infty\ \wedge\ \lim\limits_{x\to5}[f(x)\cdot g(x)]=13 \)}{\( \lim\limits_{x\to5}f(x)=\infty\ \wedge\ \lim\limits_{x\to5}g(x)=\infty\ \wedge\ \lim\limits_{x\to5}[f(x)\cdot g(x)]=13 \)}{\( \lim\limits_{x\to5}f(x)=-\infty\ \wedge\ \lim\limits_{x\to5}g(x)=\infty\ \wedge\ \lim\limits_{x\to5}[f(x)\cdot g(x)]=13 \)}{}{}}
\End

\Data\ID{1003164304}
\Updated{2020-07-15 03:07:21}
\Level{C}
\SubArea{Limits and continuity}
\LANGen{
\Question{
Which of the following situations could arise for suitable functions \( f \) and \( g \)?}
{\( \lim\limits_{x\to2} f(x)=\infty\ \wedge\ \lim\limits_{x\to2}g(x)=-\infty\ \wedge\ \lim\limits_{x\to2}[f(x)+g(x)]=-\infty \)}{\( \lim\limits_{x\to2} f(x)=13\ \wedge\ \lim\limits_{x\to2}g(x)=0\ \wedge\ \lim\limits_{x\to2}\frac{f(x)}{g(x)}=13 \)}{\( \lim\limits_{x\to2} f(x)=-\infty\ \wedge\ \lim\limits_{x\to2}g(x)=\infty\ \wedge\ \lim\limits_{x\to2}[f(x)-g(x)]=0 \)}{\( \lim\limits_{x\to2} f(x)=\infty\ \wedge\ \lim\limits_{x\to2}g(x)=-\infty\ \wedge\ \lim\limits_{x\to2}[f(x)\cdot g(x)]=\infty \)}{}{}}
\LANGpl{
\Question{
Które z poniższych sytuacji mogą zachodzić dla odpowiednich funkcji \( f \) i \( g \)?}
{\( \lim\limits_{x\to2} f(x)=\infty\ \wedge\ \lim\limits_{x\to2}g(x)=-\infty\ \wedge\ \lim\limits_{x\to2}[f(x)+g(x)]=-\infty \)}{\( \lim\limits_{x\to2} f(x)=13\ \wedge\ \lim\limits_{x\to2}g(x)=0\ \wedge\ \lim\limits_{x\to2}\frac{f(x)}{g(x)}=13 \)}{\( \lim\limits_{x\to2} f(x)=-\infty\ \wedge\ \lim\limits_{x\to2}g(x)=\infty\ \wedge\ \lim\limits_{x\to2}[f(x)-g(x)]=0 \)}{\( \lim\limits_{x\to2} f(x)=\infty\ \wedge\ \lim\limits_{x\to2}g(x)=-\infty\ \wedge\ \lim\limits_{x\to2}[f(x)\cdot g(x)]=\infty \)}{}{}}
\LANGcs{
\Question{
Která z následujících situací může pro vhodné funkce \( f \) a \( g \) nastat?}
{\( \lim\limits_{x\to2} f(x)=\infty\ \wedge\ \lim\limits_{x\to2}g(x)=-\infty\ \wedge\ \lim\limits_{x\to2}[f(x)+g(x)]=-\infty \)}{\( \lim\limits_{x\to2} f(x)=13\ \wedge\ \lim\limits_{x\to2}g(x)=0\ \wedge\ \lim\limits_{x\to2}\frac{f(x)}{g(x)}=13 \)}{\( \lim\limits_{x\to2} f(x)=-\infty\ \wedge\ \lim\limits_{x\to2}g(x)=\infty\ \wedge\ \lim\limits_{x\to2}[f(x)-g(x)]=0 \)}{\( \lim\limits_{x\to2} f(x)=\infty\ \wedge\ \lim\limits_{x\to2}g(x)=-\infty\ \wedge\ \lim\limits_{x\to2}[f(x)\cdot g(x)]=\infty \)}{}{}}
\LANGsk{
\Question{
Ktorá z nasledujúcich situácií môže pri vhodnej funkcii \( f \) a \( g \) nastať?}
{\( \lim\limits_{x\to2} f(x)=\infty\ \wedge\ \lim\limits_{x\to2}g(x)=-\infty\ \wedge\ \lim\limits_{x\to2}[f(x)+g(x)]=-\infty \)}{\( \lim\limits_{x\to2} f(x)=13\ \wedge\ \lim\limits_{x\to2}g(x)=0\ \wedge\ \lim\limits_{x\to2}\frac{f(x)}{g(x)}=13 \)}{\( \lim\limits_{x\to2} f(x)=-\infty\ \wedge\ \lim\limits_{x\to2}g(x)=\infty\ \wedge\ \lim\limits_{x\to2}[f(x)-g(x)]=0 \)}{\( \lim\limits_{x\to2} f(x)=\infty\ \wedge\ \lim\limits_{x\to2}g(x)=-\infty\ \wedge\ \lim\limits_{x\to2}[f(x)\cdot g(x)]=\infty \)}{}{}}
\LANGes{
\Question{
¿Cuál de las siguientes situaciones podría surgir para las funciones adecuadas \( f \) y \( g \)?}
{\( \lim\limits_{x\to2} f(x)=\infty\ \wedge\ \lim\limits_{x\to2}g(x)=-\infty\ \wedge\ \lim\limits_{x\to2}[f(x)+g(x)]=-\infty \)}{\( \lim\limits_{x\to2} f(x)=13\ \wedge\ \lim\limits_{x\to2}g(x)=0\ \wedge\ \lim\limits_{x\to2}\frac{f(x)}{g(x)}=13 \)}{\( \lim\limits_{x\to2} f(x)=-\infty\ \wedge\ \lim\limits_{x\to2}g(x)=\infty\ \wedge\ \lim\limits_{x\to2}[f(x)-g(x)]=0 \)}{\( \lim\limits_{x\to2} f(x)=\infty\ \wedge\ \lim\limits_{x\to2}g(x)=-\infty\ \wedge\ \lim\limits_{x\to2}[f(x)\cdot g(x)]=\infty \)}{}{}}
\End
